\documentclass[UTF8,12pt,a4paper,fontset=fandol]{ctexbook}

% 支持从项目根目录或当前笔记目录编译。
\makeatletter
\def\input@path{{../}}
\makeatother
\usepackage{researchnotes}

\title{初等数论笔记}
\author{}
\date{}

\begin{document}

\maketitle
\frontmatter
\tableofcontents

\chapter{绪言与课程路线}
\label{ch:course-outline}

\keyterm{初等数论（elementary number theory）}研究整数的整除、分解、同余及整数方程。本书面向本科生，以中学代数和基本逻辑为起点，不预设抽象代数。内容围绕三个问题展开：整数如何分解，取余运算保留哪些结构，以及整数方程的解如何判定与构造。

全书分为四个阶段。第\ref{ch:introduction}至\ref{ch:primes}章建立证明方法、整除理论与唯一分解；第\ref{ch:congruences}至\ref{ch:quadratic-residues}章依次讨论同余、中国剩余定理、算术函数、幂同余、原根与二次互反律；第\ref{ch:diophantine}至\ref{ch:pell}章讨论整数方程、连分数和 Pell 方程；第\ref{ch:prime-distribution}、\ref{ch:algorithms}章选讲素数分布与算法应用。前一阶段建立的工具会在后一阶段反复使用：贝祖等式产生模逆元，模逆元支撑中国剩余定理，素因子分解支撑算术函数，乘法阶和平方剩余又把同余方法带入整数方程。

正文突出定义、定理、证明与例题。核心推导尽量自足；较长的拓展定理在引用处明确标注不证，不以证明思路代替完整证明。章末习题用于检验概念、训练证明，附录提供代表题的答案与综合解答。计算结果用于核对例子，不能代替一般命题的证明。

学习时可以分两遍阅读。第一遍掌握每个定义的对象与条件，亲手完成例题，理解主要定理解决什么问题；第二遍追踪证明中每一步的依据，尤其注意互素、素数、非平方等条件在哪里使用。原根分类、二次互反律的格点证明、平方根连分数的周期证明较长，初读时可先把握结论和应用，再返回细读。带有“选学”的内容不作为基础习题的默认前提；素数分布中涉及级数和极限的部分需要微积分基础，其余主线以整数运算和初等代数为主。

数论题目的答案通常有几个层次。“找到一组解”只是存在性证据；“写出一族解”还要检查参数条件；“求出全部解”则必须证明任何解都能落入这族参数。阅读例题时，应分别辨认构造与完备性论证。各章习题前四题用于巩固本章主线，后续题按概念辨析、方法迁移和综合证明逐步展开；附录给出代表题的答案与后四题的关键步骤，宜在独立尝试后使用。

课程范围可与 MIT 本科课程资料\cite{mit-syllabus,mit-notes}对照，计算专题可进一步阅读 Stein 的教材\cite{stein-book}。参考资料中的代数和椭圆曲线部分不作为本书先修要求。

\mainmatter

\chapter{绪论与整数证明方法}
\label{ch:introduction}

初等数论中的计算对象很熟悉，论证方式却常与代数运算题不同。检验一百个例子仍不能证明“对所有整数都成立”，而一个反例就足以否定全称命题。本章从整数的离散性出发，说明归纳、反证、构造和抽屉原理各适合解决什么问题，为后面的证明建立共同语言。

\section{从三个问题认识数论}
\label{sec:intro-questions}

每隔 $6$ 天与每隔 $8$ 天举行的两项活动，若今天同时举行，则下一次同时举行在 $24$ 天后。这是公倍数问题。若两项活动的起始时间不同，就需要同时满足两个余数条件，这将引出同余方程。

整数 $65$ 有表示 $65=1^2+8^2=4^2+7^2$，而 $21$ 不能表示为两个整数的平方和。分数 $1/7=0.\overline{142857}$ 的循环节长度是 $6$。这些现象分别引出平方和判据和乘法阶。数论不仅寻找一个例子，更要回答：何时存在，是否唯一，怎样求出全部对象，以及计算过程为什么正确。

\section{整数、逻辑语言与符号约定}

记 $\mathbb Z$、$\mathbb Q$、$\mathbb R$ 分别为整数、有理数、实数集，$\mathbb Z_{>0}$ 和 $\mathbb Z_{\ge0}$ 分别为正整数集和非负整数集。后文还使用复数集 $\mathbb C$。字母的范围由所在命题规定，通常用 $p,q$ 表示素数，但约定不代替假设。有限集合 $A$ 的元素个数记为 $\#A$。

命题 $P\Rightarrow Q$ 表示 $P$ 是 $Q$ 的充分条件，$Q$ 是 $P$ 的必要条件；$P\Leftrightarrow Q$ 要求两个方向都成立。逆命题 $Q\Rightarrow P$ 一般不成立，逆否命题 $\neg Q\Rightarrow\neg P$ 则与原命题等价。例如，被 $4$ 整除的整数一定是偶数，但偶数 $6$ 不能被 $4$ 整除。

量词顺序不可随意交换。命题“对每个整数 $a$，存在整数 $b$ 使 $b>a$”成立，但不存在一个整数 $b$ 大于所有整数。否定全称命题只需一个反例；否定存在命题则须排除全部可能。

例如，“每个偶数都是两个奇数的和”可写为：对任意 $n\in\mathbb Z$，存在奇整数 $u,v$，使 $2n=u+v$。取 $u=1,v=2n-1$ 就完成了构造；若额外要求 $u,v$ 为正数，就必须另加 $n\ge1$，因为负偶数不可能是两个正数的和。看似细小的取值范围会改变命题本身，所以应先读清对象，再开始运算。

后文的符号 $\sum$ 和 $\prod$ 分别表示求和与求积。例如 $\sum_{j=1}^r a_j=a_1+\cdots+a_r$，$\prod_{j=1}^r a_j=a_1\cdots a_r$；空和约定为 $0$，空积约定为 $1$。这些约定使没有素因子或没有待加项的边界情形也能使用统一公式。

\section{良序原理与数学归纳法}

\begin{axiom}[良序原理]
每个非空的非负整数子集都有最小元。
\end{axiom}

\keyterm{良序原理（well-ordering principle）}是本书对整数所使用的基本性质。它既用于选择最小对象，也用于证明下降过程必定终止。

\begin{theorem}[强归纳法（strong induction）]
设 $n_0\in\mathbb Z_{\ge0}$。若命题 $P(n_0)$ 成立，且对每个 $n>n_0$，由所有 $n_0\le k<n$ 的 $P(k)$ 成立可推出 $P(n)$，则 $P(n)$ 对所有 $n\ge n_0$ 成立。
\end{theorem}
\begin{proof}
若有反例，取最小反例 $m$。由初始情形知 $m>n_0$，由最小性知 $P(k)$ 对所有 $n_0\le k<m$ 成立；归纳条件于是给出 $P(m)$，矛盾。
\end{proof}

普通归纳法只使用前一个命题 $P(n-1)$，是上述方法的特例。最小反例法把同一论证反过来组织：假定反例存在，再从最小反例构造更小的反例。任何严格递减的非负整数序列都只能有有限项，否则其值集有最小元，而下一项更小，发生矛盾。

归纳假设是一个临时前提，用于证明从已知情形走向下一情形的逻辑蕴含；它不等于预先假定整个结论成立。若构造 $n$ 的对象时只需使用 $n-1$ 的对象，普通归纳通常足够；若需要把 $n$ 拆成两个都小于 $n$、但未必等于 $n-1$ 的数，强归纳往往更自然。第\ref{ch:primes}章的素因子分解正属于后一种情形。

\begin{example}
证明对每个 $n\ge1$，$1+3+\cdots+(2n-1)=n^2$。
\end{example}
\begin{proof}
$n=1$ 时成立。若前 $n$ 项之和为 $n^2$，则前 $n+1$ 项之和为 $n^2+(2n+1)=(n+1)^2$，归纳完成。
\end{proof}

\section{反证、构造与抽屉原理}

\begin{example}
$\sqrt2$ 不是有理数。
\end{example}
\begin{proof}
若 $\sqrt2=a/b$，可在所有 $a,b>0$ 的整数表示中选取分母 $b$ 最小的一组。由 $a^2=2b^2$ 知 $a$ 为偶数：奇数 $2k+1$ 的平方仍为奇数。写 $a=2c$，得 $b^2=2c^2$，故 $b=2d$。于是 $\sqrt2=c/d$ 且 $0<d<b$，与选择矛盾。
\end{proof}

该证明同时说明了反证与递降的作用。构造性证明则必须实际给出对象，并验证它满足条件；仅说“可以找到”不构成存在性证明。

\begin{proposition}[抽屉原理]
若 $N>k\ge1$，把 $N$ 个对象放入 $k$ 个盒子，则至少一个盒子含两个对象。
\end{proposition}
\begin{proof}
若每盒至多一个，对象总数至多为 $k$，与 $N>k$ 矛盾。
\end{proof}

\keyterm{抽屉原理（pigeonhole principle）}在数论中常以余数为盒子。两个对象落在同一盒子，意味着它们的差被某个整数整除；这一思想将在同余与平方和问题中反复出现。

\begin{example}[从部分和中寻找整除关系]
给定任意 $n$ 个整数 $a_1,\ldots,a_n$，其中 $n\ge1$。证明存在一段连续项，其和被 $n$ 整除。
\end{example}
\begin{proof}
考虑 $n+1$ 个部分和 $S_0=0$、$S_j=a_1+\cdots+a_j$（$1\le j\le n$）。整数除以 $n$ 只有 $n$ 种余数，因此有 $0\le i<j\le n$ 使 $S_i,S_j$ 余数相同。于是 $S_j-S_i=a_{i+1}+\cdots+a_j$ 被 $n$ 整除。这里加入 $S_0$，使从第一项开始的连续段也包含在统一论证中。
\end{proof}

\section{从观察到证明}

发现规律与证明规律是两个相互配合的步骤。例如，前几个数 $n(n+1)$ 都是偶数，这提示一个全称命题。证明只需将 $n$ 按奇偶分类：若 $n$ 偶，第一因子被 $2$ 整除；若 $n$ 奇，第二因子被 $2$ 整除。结论依赖“两个相邻整数恰有一个偶数”，并不依赖已经计算过多少个例子。

\begin{example}[计算很多项仍不能替代证明]
有人从若干计算猜测 $n^2+n+41$ 对所有非负整数 $n$ 都是素数。怎样检验这一猜测？
\end{example}
\begin{solve}
先尝试让前两项被 $41$ 整除。取 $n=41$，得到 $41^2+41+41=41\cdot43$，已经否定猜测；取 $n=40$ 还可得到 $40^2+40+41=41^2$。反例应同时给出符合范围的输入和失败原因，而不是仅说“继续算下去可能不成立”。
\end{solve}

书写证明时可以先在草稿中倒推目标，再在正文中按逻辑依赖正向排列。证明“若 $P$ 则 $Q$”时，每一步应来自 $P$、已知定理或前一步；若使用目标 $Q$ 参与推导，应明确它是在反证假设中出现，或说明所做变形均可逆。后文将经常使用这样的组织方式：先说明要构造什么，再说明构造为什么可行，最后验证结果。

\section{习题}

本章的共同出发点是整数的离散性：最小元支持归纳与递降，有限分类支持抽屉原理。做题时先辨认量词和对象范围，再选择证明方法；计算观察可以提出猜想，不能代替对所有对象的论证。

\begin{enumerate}
  \item 写出“对每个正整数 $n$，存在素数 $p>n$”的否定，说明否定命题需要怎样的证据。
  \item 用归纳法证明 $\sum_{k=1}^n k=n(n+1)/2$，指出归纳起点与归纳假设。
  \item 证明从 $1,2,\ldots,2n$ 中任取 $n+1$ 个不同整数，必有两个相邻。
  \item 仿照最小分母论证证明 $\sqrt3$ 无理；先说明平方被 $3$ 整除时，原整数也被 $3$ 整除。
 \item\label{ex:intro-negation} 写出“存在整数 $n$，使每个整数 $m$ 都满足 $m\le n$”的否定，并证明否定命题成立。
 \item\label{ex:intro-strong} 用强归纳法证明每个整数 $n\ge2$ 都能写成 $2a+3b$，其中 $a,b$ 为非负整数；说明为什么需要两个初始情形。
 \item\label{ex:intro-dividing-pair} 从 $1,2,\ldots,2n$ 中任取 $n+1$ 个不同整数，证明有两个数满足其中一个整除另一个。提示：不断除去因子 $2$。
 \item\label{ex:intro-descent-domain} 为什么正实数序列 $1,1/2,1/4,\ldots$ 可以无限下降，而正整数序列不能？指出递降法证明必须验证的范围条件。
\end{enumerate}

\chapter{整除、最大公因数与欧几里得算法}
\label{ch:divisibility}

若有 $84$ 支笔和 $30$ 本练习本，要分成若干份完全相同的礼包且不剩余，礼包份数必须同时整除 $84$ 和 $30$；最多能分多少份，就是最大公因数问题。本章把这样的分组问题推广为整数之间的整除关系，并说明一个核心事实：求最大公因数、构造整数线性组合、求解一次整数方程，本质上使用的是同一套工具。

\section{整除关系与带余除法}

\begin{definition}
对 $a,b\in\mathbb Z$，若存在 $k\in\mathbb Z$ 使 $b=ak$，则称 $a$ \keyterm{整除（divide）} $b$，记为 $a\mid b$；否则记为 $a\nmid b$。
\end{definition}

每个整数都整除 $0$，而 $0\mid b$ 当且仅当 $b=0$。若 $a\mid b$ 且 $a\mid c$，则对任意 $u,v\in\mathbb Z$ 有 $a\mid ub+vc$，因为 $b=ar,c=as$ 蕴含 $ub+vc=a(ur+vs)$。整除具有传递性；若 $a\mid b$ 且 $b\ne0$，则 $|a|\le|b|$。

整除符号表达一种关系，不是一个商：$3\mid12$ 是真命题，而 $12/3=4$ 是数值等式。对非零 $a,b$，$a\mid b$ 且 $b\mid a$ 只能推出 $a=\pm b$，若另有 $a,b>0$ 才能推出相等。例如 $-3$ 和 $3$ 相互整除。这些符号和正负号上的区别会影响后文的唯一性陈述。

\begin{theorem}[带余除法（division algorithm）]
\label{thm:division}
对任意 $a\in\mathbb Z$ 与 $b\in\mathbb Z_{>0}$，存在唯一的整数 $q,r$，使
\begin{equation}
  a=bq+r,\qquad 0\le r<b.
  \label{eq:division-algorithm}
\end{equation}
\end{theorem}
\begin{proof}
集合 $S=\{a-bq:q\in\mathbb Z,\ a-bq\ge0\}$ 非空：例如取 $q=-|a|$。由良序原理取其最小元 $r=a-bq$。若 $r\ge b$，则 $r-b\in S$ 且更小，矛盾，故 $0\le r<b$。若另有 $a=bq'+r'$，则 $b(q-q')=r'-r$；右端绝对值小于 $b$，只能为 $0$，于是 $q=q',r=r'$。
\end{proof}

例如 $-17=5(-4)+3$，所以除以 $5$ 的余数是 $3$。余数的非负约定对负被除数同样适用。

式\eqref{eq:division-algorithm}中的 $q$ 称为商，$r$ 称为余数。若只要求 $a=bq+r$ 而不限制余数范围，就会出现无穷多个表示，例如 $17=5\cdot3+2=5\cdot2+7$；唯一性来自长度为 $b$ 的范围 $0\le r<b$。在数轴上，可以把它理解为先找到包含 $a$ 的区间 $[bq,b(q+1))$，再测量 $a$ 到左端点的整数距离。

\section{最大公因数、互素与贝祖等式}

\begin{definition}
不全为零的整数 $a,b$ 的\keyterm{最大公因数（greatest common divisor）} $\gcd(a,b)$ 是它们最大的正公因数。若 $\gcd(a,b)=1$，称 $a,b$ \keyterm{互素（coprime）}。
\end{definition}

公因数集合非空，且被任一非零数的绝对值所界，故最大公因数存在。特别地，$\gcd(a,0)=|a|$（$a\ne0$），改变符号不影响最大公因数。本书不使用 $\gcd(0,0)$。

\begin{theorem}[贝祖等式]
\label{thm:bezout}
设 $a,b\in\mathbb Z$ 不全为零，$d=\gcd(a,b)$。则存在 $u,v\in\mathbb Z$ 使 $au+bv=d$，而且每个公因数都整除 $d$。
\end{theorem}
\begin{proof}
所有正的整数线性组合 $ax+by$ 构成非空集合，取其中最小元 $s=au+bv$。将 $a$ 除以 $s$，余数 $r=a-qs$ 仍是 $a,b$ 的整数线性组合；若 $r>0$，则 $r<s$ 违背最小性，故 $s\mid a$。同理 $s\mid b$。任意公因数都整除 $s$，从而不超过 $s$；因此 $s=d$，并得到所需等式。
\end{proof}

\keyterm{贝祖等式（Bézout's identity）}把公因数问题转化为线性组合。由它还知全部整数线性组合恰为 $d$ 的整数倍。

这一证明的关键并不是直接寻找“最大的因数”，而是寻找“最小的正线性组合”。带余除法把任何不能整除的情形变成更小的正线性组合，迫使最小者同时整除 $a,b$。例如 $\gcd(12,18)=6$，所以 $12x+18y$ 的所有可能取值恰为 $6$ 的整数倍：$6=18-12$，把此式乘以任意整数即可得到每一个这样的倍数。特别地，互素等价于能用整数系数拼出 $1$。

\begin{corollary}[互素消去]
\label{cor:coprime-cancellation}
若 $\gcd(a,b)=1$ 且 $a\mid bc$，则 $a\mid c$。
\end{corollary}
\begin{proof}
取 $au+bv=1$，乘以 $c$ 得 $c=auc+bvc$；右侧两项均被 $a$ 整除。
\end{proof}

\section{欧几里得算法与扩展算法}

贝祖等式给出存在性，但其证明未直接提供便于手算的系数。欧几里得算法一方面求最大公因数，另一方面保存每次带余除法的信息，回代后就能把存在性证明变成具体计算。算法正确性要分为两件事：每次更新不改变答案，且更新一定会结束。

若 $a=bq+r$，则 $a,b$ 的公因数恰是 $b,r$ 的公因数，因而
\[
  \gcd(a,b)=\gcd(b,r).
\]
\keyterm{欧几里得算法（Euclidean algorithm）}对 $A=|a|,B=|b|$ 反复执行：若 $B=0$，输出 $A$；否则用带余除法求 $R$，令 $(A,B)\leftarrow(B,R)$。非零余数严格递减，故算法终止；最大公因数在更新中不变，最后一个非零余数即为答案。

扩展算法同时记录系数。对非负初值 $A,B$，以 $(r_0,s_0,t_0)=(A,1,0)$、$(r_1,s_1,t_1)=(B,0,1)$ 开始；当 $r_i>0$ 时取 $q_i=\lfloor r_{i-1}/r_i\rfloor$，并令
\[
 (r_{i+1},s_{i+1},t_{i+1})
 =(r_{i-1},s_{i-1},t_{i-1})-q_i(r_i,s_i,t_i).
\]
这里 $\lfloor x\rfloor$ 表示不超过实数 $x$ 的最大整数，称为下取整；后文的 $\lceil x\rceil$ 表示不小于 $x$ 的最小整数，称为上取整。不变式 $r_i=s_iA+t_iB$ 由递推保持，故最后一组非零余数的系数给出贝祖等式。原数为负时，把相应系数变号；若初始 $B=0$，直接返回第一组。

\begin{example}
求 $252,198$ 的最大公因数和一组贝祖系数。
\end{example}
\begin{solve}
带余除法依次为 $252=198+54$、$198=3\cdot54+36$、$54=36+18$、$36=2\cdot18$，故最大公因数为 $18$。回代得
\[
 18=54-36=4\cdot54-198=4\cdot252-5\cdot198.
\]
\end{solve}

同一计算也可用系数表完成。下表的每一行都满足 $r=252s+198t$，下一行由前两行相减若干倍得到；因此读到最后一个非零余数时，贝祖系数已经同时求出。
\[
\begin{array}{c|r|r|r}
 &r&s&t\\\hline
0&252&1&0\\
1&198&0&1\\
2&54&1&-1\\
3&36&-3&4\\
4&18&4&-5
\end{array}
\]
回代法较适合短计算，系数表适合写成程序。二者记录的是同一个不变式，并不是两个不同的定理。

\section{最小公倍数与一次不定方程}

非零整数 $a,b$ 的\keyterm{最小公倍数（least common multiple）}记为 $\operatorname{lcm}(a,b)$，即最小正公倍数。写 $|a|=dA,|b|=dB$，其中 $d=\gcd(a,b)$，则 $\gcd(A,B)=1$。若 $m$ 为公倍数，写 $m=dAt$，由 $B\mid At$ 得 $B\mid t$，故 $dAB\mid m$。因此
\begin{equation}
 \operatorname{lcm}(a,b)=\frac{|ab|}{\gcd(a,b)}.
 \label{eq:gcd-lcm}
\end{equation}

多个整数的最大公因数可递归计算；它们整体最大公因数为 $1$ 并不等于两两互素，例如 $6,10,15$。

\begin{theorem}[一次不定方程]
\label{thm:linear-diophantine}
设 $a,b,c\in\mathbb Z$，$a,b$ 不全为零，$d=\gcd(a,b)$。方程 $ax+by=c$ 有整数解当且仅当 $d\mid c$。若 $(x_0,y_0)$ 是一组解，则全部解为
\begin{equation}
 x=x_0+\frac bd t,\qquad y=y_0-\frac ad t,\qquad t\in\mathbb Z.
 \label{eq:linear-diophantine-solutions}
\end{equation}
\end{theorem}
\begin{proof}
必要性由 $d\mid ax+by$ 得到；充分性将贝祖等式乘以 $c/d$ 即可。任意两组解之差满足 $(a/d)(x-x_0)=-(b/d)(y-y_0)$。若 $b\ne0$，互素消去给出 $b/d\mid x-x_0$，代入得到通解；若 $b=0$，则 $a/d=\pm1$，$x$ 固定而 $y$ 任意，同一公式仍覆盖全部解。反向代入验证每组参数确为解。
\end{proof}

要求未知量取整数值的方程称为\keyterm{不定方程（Diophantine equation）}；名称不意味着它必有无穷多解。

通解中的参数 $t$ 有明确的几何含义：从一组整数解出发，沿着向量 $(b/d,-a/d)$ 走整数步，仍在直线 $ax+by=c$ 上。除以 $d$ 后两个方向分量互素，保证这个步长没有跳过其他整数点。求非负解时，应先得到整数通解，再把 $x,y\ge0$ 转化为对同一个参数 $t$ 的上下界；分别寻找两个看似合适的参数会破坏两坐标之间的关联。

\begin{example}
求 $18x+30y=42$ 的全部整数解和非负整数解。
\end{example}
\begin{solve}
约去最大公因数 $6$ 后得到 $3x+5y=7$，一组解为 $(-1,2)$，故 $x=-1+5t,y=2-3t$。非负性要求 $t\ge1$ 且 $t\le0$，所以不存在非负整数解，尽管整数解有无穷多组。
\end{solve}

\section{非负表示与两种面额问题}

在计件和配比问题中，负数解通常没有实际含义。例如，用面额 $4$ 与 $7$ 的单位凑出总量 $N$，要求的是 $4x+7y=N$ 且 $x,y\ge0$。互素只保证整数解存在，并不保证非负解存在；不过当总量足够大时，情况会稳定下来。

\begin{theorem}[两种面额的最大不可表示数]
\label{thm:two-coins}
若整数 $a,b\ge2$ 且 $\gcd(a,b)=1$，则不能表示成 $ax+by$（$x,y\in\mathbb Z_{\ge0}$）的最大整数为 $ab-a-b$。换言之，$ab-a-b$ 不可表示，而每个整数 $N>ab-a-b$ 都可表示。
\end{theorem}
\begin{proof}
若 $ab-a-b=ax+by$，移项得 $ab=a(x+1)+b(y+1)$。于是 $a\mid b(y+1)$，互素消去给出 $a\mid y+1$。因为 $y+1>0$，有 $y+1\ge a$，从而右侧严格大于 $ab$，矛盾。

再设 $N>ab-a-b$。数 $0,b,2b,\ldots,(a-1)b$ 除以 $a$ 的余数两两不同：若 $a\mid b(i-j)$，则 $a\mid i-j$，而 $|i-j|<a$，故 $i=j$。它们恰好覆盖全部余数，因此可选 $0\le y\le a-1$ 使 $a\mid N-by$。写 $N-by=ax$。由
\[
 N-by\ge N-(a-1)b>-a
\]
知 $ax$ 是严格大于 $-a$ 的 $a$ 的整数倍，只能满足 $ax\ge0$。于是 $x,y$ 均非负。
\end{proof}

\begin{example}[存在性判据与实际求解]
用 $4$ 与 $7$ 两种单位表示 $25$，求全部非负表示；再判断哪些较大的总量一定可表示。
\end{example}
\begin{solve}
一组整数解为 $(1,3)$，由一次不定方程的通解公式，可写 $x=1-7t$、$y=3+4t$，其中 $t\in\mathbb Z$。非负性迫使 $t=0$，唯一表示为 $25=4+3\cdot7$。定理\ref{thm:two-coins}又说明最大不可表示数为 $28-4-7=17$，因而每个至少为 $18$ 的整数都可表示。定理判定大范围的存在性，通解则回答某个总量究竟有多少种表示。
\end{solve}

\section{习题}

本章把公因数、线性组合和一次整数方程联系起来。欧几里得算法提供具体系数，贝祖等式解释可解性；非负条件则需要在整数通解上另行限制参数。解题时应把“有整数解”与“有符合实际范围的解”分开。

\begin{enumerate}
 \item 对 $-37$ 除以 $8$ 作带余除法；求 $\gcd(414,662)$ 并写出贝祖等式。
 \item 证明 $\gcd(a,b)=\gcd(a,a+b)$，其中 $a,b$ 不全为零。
 \item 求 $35x+22y=1$ 的全部整数解。
 \item 若 $a,b>0$ 且互素，证明 $a\mid c$、$b\mid c$ 蕴含 $ab\mid c$。
 \item\label{ex:div-consecutive} 对任意整数 $n$，证明 $\gcd(n,n+1)=1$，并直接给出一组贝祖系数。
 \item\label{ex:div-combinations} 对任意整数 $n$，证明 $\gcd(2n+1,3n+1)=1$。提示：寻找与 $n$ 无关的线性组合。
 \item\label{ex:div-nonnegative} 求 $4x+7y=100$ 的全部非负整数解，并解释为什么不会遗漏。
 \item\label{ex:div-gcd-lcm-count} 给定正整数 $d,L$，求满足 $\gcd(a,b)=d$、$\operatorname{lcm}(a,b)=L$ 的正整数有序对 $(a,b)$ 的个数。可使用下一章的唯一分解，并按 $d\mid L$ 是否成立分类。
\end{enumerate}

\chapter{素数与算术基本定理}
\label{ch:primes}

正整数可以反复拆成较小整数的乘积，直到不能继续拆分。这样的终点就是素数。本章要证明两件不同的事：分解总能进行到底，而且无论怎样分解，得到的素因子及其重数都相同。唯一分解使许多关于整数的命题转化为关于非负整数指数的命题，是后续整除计算的共同基础。

\section{素数、合数与素数无穷性}

\begin{definition}
大于 $1$ 且正因数只有 $1$ 和自身的整数称为\keyterm{素数（prime number）}；大于 $1$ 的其他整数称为合数（composite number）。整数 $1$ 不属于这两类。
\end{definition}

每个 $n>1$ 都有素因子：取 $n$ 的大于 $1$ 的最小正因数 $d$；若 $d$ 为合数，它有一个介于 $1$ 与 $d$ 之间的因数，这也是 $n$ 的因数，矛盾。

\begin{theorem}[素数无穷性]
素数有无穷多个。
\end{theorem}
\begin{proof}
若全部素数为 $p_1,\ldots,p_r$，令 $N=p_1\cdots p_r+1>1$。取 $N$ 的素因子 $q$，则 $q$ 应是某个 $p_i$，于是同时整除 $p_1\cdots p_r$ 和 $N$，从而整除 $1$，矛盾。
\end{proof}

这里并未断言 $N$ 本身为素数，只断言它具有新的素因子。

例如 $2\cdot3\cdot5\cdot7\cdot11\cdot13+1=30031=59\cdot509$。它虽是合数，两个因子却都不在原先的列表中，足以实现证明目的。理解一个构造时，应检查证明究竟需要“新数是素数”，还是只需要“新数有新的素因子”；前一个要求在这里既不成立，也无必要。

\section{欧几里得引理与唯一分解}

\begin{lemma}[欧几里得引理（Euclid's lemma）]
\label{lem:euclid-prime}
若 $p$ 为素数且 $p\mid ab$，则 $p\mid a$ 或 $p\mid b$。
\end{lemma}
\begin{proof}
若 $p\nmid a$，由素数定义有 $\gcd(p,a)=1$，应用推论\ref{cor:coprime-cancellation}得 $p\mid b$。
\end{proof}

\begin{theorem}[算术基本定理（fundamental theorem of arithmetic）]
\label{thm:unique-factorization}
每个整数 $n>1$ 都有唯一表示
\begin{equation}
 n=p_1^{\alpha_1}\cdots p_r^{\alpha_r},\qquad
 p_1<\cdots<p_r,\quad \alpha_i\in\mathbb Z_{>0},
 \label{eq:prime-factorization}
\end{equation}
其中 $p_1,\ldots,p_r$ 为素数。
\end{theorem}
\begin{proof}
对 $n$ 强归纳证明存在性。若 $n$ 为素数，已成立；否则 $n=ab$，其中 $1<a,b<n$，把 $a,b$ 的素数分解相乘即可。为证唯一性，将任意两种分解展开为素数的乘积。第一种的任一素因子，由欧几里得引理必整除第二种中的某个素数，因而等于它；消去这一对因子后继续，最终所有素因子及重数一致。
\end{proof}

对负整数，另加一个负号即可。整数 $1$ 对应空乘积，值约定为 $1$。把 $1$ 排除在素数之外，正是为了使分解中的因子个数不受任意插入 $1$ 的影响。

唯一性证明中最关键的是素数的消去性质，不能把它直接推广给任意整数。例如 $6\mid2\cdot3$，但 $6$ 既不整除 $2$，也不整除 $3$。此外，两种分解不可能在消去公共因子后出现“一边还有素数，另一边为 $1$”：若一边还有因子，其乘积至少为 $2$。这一点补足了逐次消去论证的终止情形。

\section{素因子指数与整除判定}

对素数 $p$ 和非零整数 $n$，记 $v_p(n)$ 为 $p$ 在 $|n|$ 的分解中的指数，称为 $p$ 进赋值（$p$-adic valuation）。不出现的素因子指数为 $0$。由唯一分解，非零 $a,b$ 满足
\[
 v_p(ab)=v_p(a)+v_p(b),\qquad
 a\mid b\ \Longleftrightarrow\ v_p(a)\le v_p(b)\ \text{对每个素数 }p.
\]
因此最大公因数与最小公倍数的指数分别取两者的最小值和最大值。

若 $a,b,a+b$ 均非零，写 $a=p^ru,b=p^sv$ 且 $p\nmid uv$。当 $r<s$ 时，$a+b=p^r(u+p^{s-r}v)$，括号内不被 $p$ 整除；故
\[
 v_p(a+b)\ge\min\{v_p(a),v_p(b)\},
\]
且两指数不同时等号成立。指数相同则可能严格增大，例如 $v_2(1+3)=2$。

\begin{example}[把整除要求转化为指数要求]
求最小正整数 $k$，使 $72k$ 为完全平方数；再求最小正整数 $t$，使 $72\mid t^2$。
\end{example}
\begin{solve}
$72=2^3\cdot3^2$。第一问要求 $3+v_2(k)$ 和 $2+v_3(k)$ 均为偶数，并且其他素因子的指数也为偶数。最小选择是 $v_2(k)=1$、其余指数为零，故 $k=2$，此时 $72k=144$。第二问要求 $2v_2(t)\ge3$、$2v_3(t)\ge2$，所以 $v_2(t)\ge2$、$v_3(t)\ge1$，最小 $t=12$。两个问题都得到平方数 $144$，但变量满足的是不同的指数条件。
\end{solve}

\begin{corollary}
固定 $k\ge2$。正整数 $n$ 是整数的 $k$ 次幂，当且仅当每个 $v_p(n)$ 都是 $k$ 的倍数。特别地，互素正整数的乘积为平方时，各因子均为平方。
\end{corollary}
\begin{proof}
若 $n=t^k$，则 $v_p(n)=kv_p(t)$；反之把各指数除以 $k$ 后构造 $t$。互素因子不共享素因子，故乘积中的偶指数分别属于各因子。
\end{proof}

\section{阶乘、二项式系数与筛法}

\begin{proposition}[阶乘中的素因子指数]
对素数 $p$ 与整数 $n\ge0$，
\begin{equation}
 v_p(n!)=\sum_{j\ge1}\left\lfloor\frac{n}{p^j}\right\rfloor.
 \label{eq:factorial-valuation}
\end{equation}
\end{proposition}
\begin{proof}
$n!$ 中每个 $p$ 的倍数至少贡献一次，每个 $p^2$ 的倍数再贡献一次，依此类推。若某数的指数为 $s$，它恰好被计数 $s$ 次。$p^j>n$ 后各项为零，故这里只涉及有限求和；$n=0$ 时两边均为零。
\end{proof}

对 $0\le k\le n$，二项式系数 $\binom nk=n!/[k!(n-k)!]$ 的素因子指数由三次应用式\eqref{eq:factorial-valuation}得到。例如 $100!$ 中 $5$ 的指数为 $20+4=24$，$2$ 的指数更大，故十进制末尾恰有 $24$ 个零。

\begin{example}[分层计数为何不会重复多算]
计算 $v_2(10!)$，并说明因子 $8$ 被计数三次的原因。
\end{example}
\begin{solve}
偶数 $2,4,6,8,10$ 各贡献一个 $2$；$4,8$ 各再贡献一个；$8$ 再贡献一个。因此总指数为 $5+2+1=8$。因子 $8=2^3$ 本来就应贡献三个 $2$，三次出现恰好记录它的三层因子。这个计数方式统计的是素因子的重数，统计对象并不是互不重复的整数。
\end{solve}

二项式系数的整性也可先从组合意义理解：从 $n$ 个不同对象中选出 $k$ 个，不考虑次序，共有 $\binom nk$ 种方法。于是阶乘商确实是整数，可以使用素因子指数判定其整除性。例如
\[
 v_2\!\left(\binom{10}{4}\right)
 =(5+2+1)-(2+1)-(3+1)=1,
\]
与 $\binom{10}{4}=210$ 相符。

若 $n$ 为合数，写 $n=ab$ 且 $1<a\le b$，则 $a\le\sqrt n$；$a$ 的一个素因子也不超过 $\sqrt n$。因此试除只需检查不超过平方根的素数。\keyterm{埃拉托色尼筛法（sieve of Eratosthenes）}用于列出 $N$ 以下的所有素数：从 $2$ 起，遇到尚未标记的 $p$ 就保留它，并标记 $p^2,p^2+p,\ldots\le N$；处理至 $p^2>N$ 后停止。较小的合数倍数已由更小素因子标记，所有合数最终都会被某个不超过其平方根的素因子筛去。

\section{习题}

唯一分解把整除关系变为各素数指数的比较：乘法对应相加，最大公因数取最小指数，最小公倍数取最大指数。使用阶乘公式时，计数对象是素因子的重数；讨论素性时，找不到小因子必须配合平方根上界才能得出结论。

\begin{enumerate}
 \item 分解 $360$；证明它的每个正约数唯一写成 $2^a3^b5^c$，其中 $0\le a\le3,0\le b\le2,0\le c\le1$。
 \item 计算 $v_2(100!)$ 与 $v_5\bigl(\binom{100}{50}\bigr)$。
 \item 若正整数 $a,b$ 互素且 $ab$ 为立方，证明 $a,b$ 都是立方。
 \item 证明对素数 $p$ 和 $1\le k<p$，$p\mid\binom pk$。
 \item\label{ex:prime-cube} 求最小正整数 $k$，使 $180k$ 为整数的立方。
 \item\label{ex:prime-odd-divisors} 证明正整数的正约数个数为奇数，当且仅当它为完全平方数。分别用约数配对与素因子指数证明。
 \item\label{ex:prime-rational-root} 证明若正整数 $n$ 的平方根为有理数，则 $n$ 为完全平方数。
 \item\label{ex:prime-binomial} 对任意素数 $p$，计算 $v_p\bigl(\binom{p^2}{p}\bigr)$，并说明结果对 $p=2$ 是否成立。
\end{enumerate}

\chapter{同余与剩余类}
\label{ch:congruences}

钟面上经过 $12$ 小时会回到原来的刻度；判断星期时，只关心天数除以 $7$ 的余数。同余把这种“忽略整圈，只保留余数”的运算写成严格的数学语言。它使无限多个整数归入有限个类，因而适合处理大数、周期以及整数方程的无解判定。不过，模运算中的除法需要额外条件，这是本章最应注意的区别。

\section{同余关系及其基本运算}

\begin{definition}
固定整数 $m\ge2$。若 $m\mid(a-b)$，称整数 $a,b$ 模 $m$ \keyterm{同余（congruent）}，记为 $a\equiv b\pmod m$。
\end{definition}

由带余除法，两个整数同余当且仅当它们除以 $m$ 的余数相同。整除的基本性质给出自反性、对称性与传递性，故同余是等价关系。需要模 $1$ 时，约定任意两个整数都同余。

例如 $17\equiv5\pmod{12}$，因为 $17-5=12$；也有 $-7\equiv5\pmod{12}$，因为 $-7-5=-12$。这里并没有说 $17=5$，而是说它们保留的余数信息相同。表达式 $a\bmod m$ 在计算语境中通常指 $a$ 的最小非负余数，是一个具体整数；$a\equiv b\pmod m$ 则是比较两个整数的命题。严格区分这两种用法，可以避免把剩余类与代表元混为一谈。

\begin{proposition}
若 $a\equiv b\pmod m$、$c\equiv d\pmod m$，则 $a\pm c\equiv b\pm d\pmod m$ 且 $ac\equiv bd\pmod m$。因而对每个整系数多项式 $f$，$f(a)\equiv f(b)\pmod m$。
\end{proposition}
\begin{proof}
和差结论直接来自差的整除性；乘积使用 $ac-bd=c(a-b)+b(c-d)$。重复使用乘法与加法即可得到多项式结论。
\end{proof}

例如 $10\equiv-1\pmod{11}$，所以 $10^{2k}\equiv1$、$10^{2k+1}\equiv-1\pmod{11}$。取余可在运算的每一步进行，不必先算出大整数。

\section{同余约分与条件的作用}

\begin{proposition}[约分法则]
\label{prop:congruence-cancellation}
设 $a,b,c\in\mathbb Z$、$m\ge2$，$d=\gcd(c,m)$。则
\begin{equation}
 ac\equiv bc\pmod m
 \quad\Longleftrightarrow\quad
 a\equiv b\pmod{m/d}.
 \label{eq:congruence-cancellation}
\end{equation}
\end{proposition}
\begin{proof}
写 $c=dc',m=dm'$，则 $\gcd(c',m')=1$。原条件等价于 $m'\mid c'(a-b)$，由互素消去又等价于 $m'\mid a-b$；$m'=1$ 时也成立。
\end{proof}

因此，只有已知 $\gcd(c,m)=1$ 时，才能普遍地约去 $c$ 而保留模数。例如 $2\cdot1\equiv2\cdot4\pmod6$，约分后得到 $1\equiv4\pmod3$，不能得到模 $6$ 的同余。另一个常用规则是：若 $n\mid m$，则模 $m$ 同余蕴含模 $n$ 同余，反向通常不成立。

\begin{example}[约分同时改变了什么]
从 $12x\equiv18\pmod{30}$ 出发求解，并解释为什么不能直接写成 $2x\equiv3\pmod{30}$。
\end{example}
\begin{solve}
原式表示 $30\mid6(2x-3)$，等价于 $5\mid2x-3$，故应写为 $2x\equiv3\pmod5$。由于 $2\cdot3\equiv1\pmod5$，两边乘以 $3$ 得 $x\equiv4\pmod5$。例如 $x=4$ 满足原式，却不满足错误的 $2x\equiv3\pmod{30}$。若约去的数与模数有公因子，除去系数中的因子时，也必须除去模数中相应的公因子。
\end{solve}

\section{剩余类、剩余系与模逆元}

\begin{definition}
整数 $a$ 的模 $m$ \keyterm{剩余类（residue class）}为 $[a]_m=\{a+km:k\in\mathbb Z\}$。全部剩余类的集合记为 $\mathbb Z/m\mathbb Z$。每类各选一个代表元所得的集合称为完全剩余系（complete residue system）。
\end{definition}

$0,1,\ldots,m-1$ 构成完全剩余系，所以恰有 $m$ 个类。定义 $[a]_m+[b]_m=[a+b]_m$、$[a]_m[b]_m=[ab]_m$；前节的保同余性质保证结果不依赖代表元。

“不依赖代表元”是定义这两种运算必须验证的条件。例如模 $5$ 有 $[2]_5=[7]_5$、$[4]_5=[-1]_5$，用前一组代表元相加得到 $[6]_5=[1]_5$，用后一组得到的也是 $[6]_5$；一般情形由同余的加法性质保证。剩余类本身是一个无穷集合，而 $2$、$7$ 只是同一集合中的不同代表。后文为简化书写，常直接用代表元表示类，但“唯一解”通常指唯一的类。

\begin{theorem}[模逆元]
\label{thm:modular-inverse}
整数 $a$ 模 $m$ 存在逆元，即存在 $u$ 使 $au\equiv1\pmod m$，当且仅当 $\gcd(a,m)=1$；逆元在模 $m$ 意义下唯一。
\end{theorem}
\begin{proof}
若 $au-mv=1$，则 $a,m$ 的每个公因数都整除 $1$，故互素。反之用贝祖等式构造 $u$。若 $au\equiv av\equiv1$，由约分法则知 $u\equiv v\pmod m$。
\end{proof}

\keyterm{模逆元（modular inverse）}记为 $a^{-1}\pmod m$，这里的负指数不是实数倒数。每个可逆类各选一个代表元，称为简化剩余系（reduced residue system）。模素数时每个非零类均可逆，模合数时则可能存在两个非零类相乘为零，如 $[2]_6[3]_6=[0]_6$。

这样的非零类称为\keyterm{零因子（zero divisor）}。它解释了为什么合数模数下不能随意约分：乘以零因子可能把不同的类变成同一类。反之，乘以可逆类可以再乘其逆元恢复原类，因此不会丢失信息。模素数时的许多性质优于模合数时，根本原因正是每个非零类均可逆。

\begin{example}
求 $3$ 模 $7$ 的逆元，并比较模 $7$ 与模 $8$ 的可逆类。
\end{example}
\begin{solve}
$3\cdot5-7\cdot2=1$，所以逆元为 $5$。模 $7$ 的 $1,\ldots,6$ 全部可逆；模 $8$ 只有 $1,3,5,7$ 可逆，因为这些数且仅这些数与 $8$ 互素。
\end{solve}

\section{数位、周期与同余不变量}

设非负整数 $N=\sum_{j=0}^k a_j10^j$，其中 $a_j$ 为十进制数位。由 $10\equiv1\pmod9$ 与 $10\equiv-1\pmod{11}$ 得
\[
 N\equiv\sum_{j=0}^k a_j\pmod9,\qquad
 N\equiv\sum_{j=0}^k(-1)^ja_j\pmod{11}.
\]
这分别给出数位和与交错数位和判别法；模 $3$ 同理。一般进制 $B$ 的展开可用相同方法处理。

\begin{example}[用不同模数检查同一个整数]
判断 $72831$ 是否被 $9$ 与 $11$ 整除。
\end{example}
\begin{solve}
数位和为 $7+2+8+3+1=21$，所以原数模 $9$ 余 $3$，不能被 $9$ 整除。从个位起作交错和，得到 $1-3+8-2+7=11$，所以原数被 $11$ 整除。相同的十进制展开在不同模数下保留不同信息；一个判别法失败，不影响另一个判别法成立。
\end{solve}

幂序列 $1,a,a^2,\ldots$ 模 $m$ 只有有限个状态，故存在 $0\le i<j$ 使 $a^i\equiv a^j$；之后乘以相同的幂，得到从第 $i$ 项起的周期性。若 $a$ 可逆，则可约去 $a^i$，得到从起点开始的周期性。一般底数只有最终周期性，例如模 $8$ 的 $1,2,4,0,0,\ldots$。

\begin{example}
证明 $x^2+y^2=4z+3$ 没有整数解。
\end{example}
\begin{proof}
偶数平方模 $4$ 余 $0$，奇数平方模 $4$ 余 $1$，所以两平方和只能余 $0,1,2$，而右侧余 $3$，矛盾。
\end{proof}

\section{习题}

同余的加、减、乘总能保持余数信息，约分却可能改变模数。可逆类允许恢复乘法前的信息，零因子则可能合并不同类。用同余排除整数解时，应明确列出所有可能余数，避免只检查几个样本。

\begin{enumerate}
 \item 求 $17$ 模 $43$ 的逆元。
 \item 证明任意 $n+1$ 个整数中有两个之差被 $n$ 整除，其中 $n\ge1$。
 \item 求整数平方模 $8$ 的全部可能余数，并证明奇数平方模 $8$ 均余 $1$。
 \item 对进制 $B\ge2$，推导模 $B-1$ 与模 $B+1$ 的数位判别法。
 \item\label{ex:cong-cancel} 求 $4x\equiv4\pmod{12}$ 的全部解类，并指出把 $4$ 直接约去但保留模数会漏掉哪些解。
 \item\label{ex:cong-obstruction} 证明 $x^2+1=3y$ 没有整数解；再判断 $x^2+1=5y$ 是否有整数解。
 \item\label{ex:cong-last-digit} 求 $7^{2026}$ 的个位数，写出所使用的余数周期。
 \item\label{ex:cong-palindrome} 证明十进制下位数为偶数的回文整数都被 $11$ 整除。回文是指从左向右与从右向左读数位相同。
\end{enumerate}

\chapter{一次同余方程与中国剩余定理}
\label{ch:crt}

设某项活动从第 $1$ 天起每隔 $6$ 天举行，另一项从第 $5$ 天起每隔 $8$ 天举行；它们在哪些天相遇？日期必须同时满足两个同余条件。与普通一次方程不同，同余方程的解通常是一个或几个周期性的整数集合。本章先解决一条线性约束，再说明多条约束怎样合并，以及相互冲突时如何识别。

\section{一次同余方程的可解性与解数}

\begin{theorem}
\label{thm:linear-congruence}
设 $a,b\in\mathbb Z$、$m\ge2$，$d=\gcd(a,m)$。方程 $ax\equiv b\pmod m$ 可解当且仅当 $d\mid b$。可解时，全部整数解构成模 $m/d$ 的一个同余类，在模 $m$ 意义下恰有 $d$ 个不同解类。
\end{theorem}
\begin{proof}
方程等价于存在整数 $y$ 使 $ax-my=b$，由定理\ref{thm:linear-diophantine}得到可解条件。约去 $d$ 后，$a/d$ 与 $m/d$ 互素，因此当 $m/d>1$ 时，模逆元给出唯一的 $x_0\pmod{m/d}$；若 $m/d=1$，任意整数都是解。模 $m$ 的代表元可取
\[
 x_0,\ x_0+m/d,\ \ldots,\ x_0+(d-1)m/d.
\]
它们两两不同，且覆盖全部整数解。
\end{proof}

例如 $6x\equiv9\pmod{15}$ 约为 $2x\equiv3\pmod5$，故 $x\equiv4\pmod5$；在模 $15$ 下有三个解类 $4,9,14$。求解时必须说明采用哪个模数计数。

从映射的角度看，模 $m$ 的每个类 $[x]_m$ 经乘以 $a$ 变为 $[ax]_m$。当 $a$ 与 $m$ 互素时，这个映射可逆，每个右端恰对应一个类；当最大公因数为 $d>1$ 时，某些右端永远无法得到，能得到的右端则各有 $d$ 个原像。定理给出的“无解或恰有 $d$ 类”正是这一现象的准确计数。例如 $6x\equiv8\pmod{15}$ 无解，因为左侧与 $15$ 的公共因子 $3$ 不整除右侧。

\section{互素模数的中国剩余定理}

\begin{theorem}[中国剩余定理]
\label{thm:crt}
设 $m_1,\ldots,m_r\ge2$ 两两互素，$a_i\in\mathbb Z$，$M=\prod_{i=1}^r m_i$。则
\begin{equation}
 x\equiv a_i\pmod{m_i},\qquad 1\le i\le r,
 \label{eq:crt-system}
\end{equation}
在模 $M$ 意义下有唯一解。
\end{theorem}
\begin{proof}
令 $M_i=M/m_i$。由两两互素，$\gcd(M_i,m_i)=1$，取 $u_i$ 使 $M_iu_i\equiv1\pmod{m_i}$。则
\begin{equation}
 x_0=\sum_{i=1}^r a_iM_iu_i
 \label{eq:crt-construction}
\end{equation}
满足每个方程，因为其他各项均被相应的 $m_i$ 整除。若 $x,y$ 都为解，每个 $m_i$ 都整除 $x-y$，两两互素使其乘积 $M$ 也整除 $x-y$，唯一性得证。
\end{proof}

\keyterm{中国剩余定理（Chinese remainder theorem）}还说明，模 $M$ 的剩余类与各个模 $m_i$ 剩余类的组合一一对应，而且这种对应保持加法和乘法。

构造式中的 $M_iu_i$ 可以看成一个“只保留第 $i$ 条信息”的系数：模 $m_i$ 时它为 $1$，模其他 $m_j$ 时它为 $0$。乘以指定余数 $a_i$ 后，再把这些项相加，各个余数互不干扰。唯一性部分则说明，虽然求逆元时可以选不同代表元、求出的整数 $x_0$ 也可能不同，但它们最后都属于同一个模 $M$ 的类。

\begin{example}
求模 $3,5,7$ 分别余 $2,3,2$ 的全部整数。
\end{example}
\begin{solve}
$M=105$，可取 $M_1u_1=35\cdot2=70$、$M_2u_2=21$、$M_3u_3=15$。式\eqref{eq:crt-construction}给出 $x_0=2\cdot70+3\cdot21+2\cdot15=233$，故全部解为 $x\equiv23\pmod{105}$。
\end{solve}

同一例题也可逐步代入。由第一式写 $x=2+3t$，代入第二式得 $3t\equiv1\pmod5$，故 $t=2+5u$，于是 $x=8+15u$。再代入模 $7$ 的条件，得 $1+u\equiv2\pmod7$，所以 $u=1+7v$，最终 $x=23+105v$。直接构造适合统一证明，逐步代入更便于小规模手算；两种方法最后都应把自由整数参数或最终模数写清楚。

\section{非互素模数的相容条件}

\begin{theorem}
\label{thm:general-crt}
方程组 $x\equiv a_i\pmod{m_i}$（$m_i\ge2$）可解，当且仅当对任意 $i,j$，
\[
 \gcd(m_i,m_j)\mid(a_i-a_j).
\]
若可解，解在模 $L=\operatorname{lcm}(m_1,\ldots,m_r)$ 意义下唯一。
\end{theorem}
\begin{proof}
任何解都给出 $a_i-a_j=(a_i-x)+(x-a_j)$，所以相容条件必要。反之，对每个整除 $L$ 的素数 $p$，选取使 $v_p(m_i)$ 最大的指标 $i(p)$，规定
\[
 x\equiv a_{i(p)}\pmod{p^{v_p(L)}}.
\]
这些素数幂两两互素，故有公共解。对任意 $m_j$ 的每个素数幂因子，相容条件保证此解与 $a_j$ 同余，于是满足全部原方程。任意两解之差被每个 $m_i$ 整除，故被 $L$ 整除。
\end{proof}

实际计算可逐项合并。已有 $x\equiv A\pmod L$，再加入 $x\equiv a\pmod m$，令 $x=A+Lt$，只需解 $Lt\equiv a-A\pmod m$。若不可解则停止；若得到 $t_0$，更新解为 $A+Lt_0$，更新模数为 $\operatorname{lcm}(L,m)$。

\begin{example}[有公共因子的两条周期约束]
求章首两项活动相遇的日期，即解 $x\equiv1\pmod6$、$x\equiv5\pmod8$。
\end{example}
\begin{solve}
先检查相容性：$\gcd(6,8)=2$，且 $2\mid5-1$，所以有解。写 $x=1+6t$，得 $6t\equiv4\pmod8$，等价于 $3t\equiv2\pmod4$，所以 $t=2+4s$，从而 $x=13+24s$。最小正日期为 $13$，之后每隔 $24$ 天重复。若把第二个起始日期改为 $4$，第一条要求 $x$ 为奇数，第二条要求 $x$ 为偶数，立即无解；这就是相容条件的具体含义。
\end{solve}

\section{素数幂分解与解的组合}

若 $m=\prod_i p_i^{\alpha_i}$，整系数多项式同余 $f(x)\equiv0\pmod m$ 等价于所有 $f(x)\equiv0\pmod{p_i^{\alpha_i}}$。局部根的每一种组合恰好给出一个模 $m$ 的根类，因此总体根数等于局部根数的乘积。

\begin{example}
求 $x^2\equiv1\pmod{15}$ 的全部解。
\end{example}
\begin{solve}
模 $3$ 与模 $5$ 分别只有 $\pm1$ 两根。将四种符号组合用中国剩余定理合并，得到 $x\equiv1,4,11,14\pmod{15}$。这也说明合数模数下二次同余可有超过两个根。
\end{solve}

这里的四种组合分别为 $(1,1),(1,-1),(-1,1),(-1,-1)$，对应 $x=1,4,11,14$。不能在分解 $(x-1)(x+1)\equiv0\pmod{15}$ 后断言某一个因子必被 $15$ 整除；例如 $x=4$ 时，$3$ 与 $5$ 分别整除不同的因子。将模数拆成互素的素数幂，正是为了分别处理这些局部选择。

\section{习题}

一次同余首先检查最大公因数是否整除右端，再求解并注明最终模数。多条约束在互素时可以任意拼合，在不互素时必须检查共享余数信息是否一致。素数幂分解则把非线性同余的解分解为独立的局部选择。

\begin{enumerate}
 \item 解 $14x\equiv8\pmod{30}$，列出模 $30$ 的全部解类。
 \item 解 $x\equiv1\pmod6$、$x\equiv5\pmod8$；再将第二式的 $5$ 改为 $4$，判断可解性。
 \item 求 $x^2\equiv1\pmod{35}$ 的全部解。
 \item 设 $p,q$ 为不同奇素数，证明 $x^2\equiv1\pmod{pq}$ 恰有四个根类。
 \item\label{ex:crt-redundant} 解 $x\equiv2\pmod4$、$x\equiv6\pmod8$，说明其中哪条条件已被另一条包含。
 \item\label{ex:crt-three} 解 $x\equiv2\pmod6$、$x\equiv5\pmod9$、$x\equiv4\pmod5$，给出最小非负解和最终模数。
 \item\label{ex:crt-zero-roots} 求 $x^2\equiv0\pmod{12}$ 的全部根类，解释为什么不能只检查 $x\equiv0\pmod{12}$。
 \item\label{ex:crt-idempotents} 求 $x^2\equiv x\pmod{30}$ 的全部根类；推广证明若 $m$ 为 $r$ 个不同素数的乘积，则 $x^2\equiv x\pmod m$ 恰有 $2^r$ 个根类。
\end{enumerate}

\chapter{算术函数与莫比乌斯反演}
\label{ch:arithmetic-functions}

前面研究一个整数能否整除另一个整数，本章转而统计一个整数附带的数量：有多少个约数、约数之和是多少、一个周期中有多少个数与它互素。唯一分解使这些数量常能拆成素数幂上的计算。莫比乌斯反演进一步回答一个反向问题：如果已知每个整数的“约数累计量”，能否恢复每个约数原本的贡献？

\section{算术函数、积性与约数求和}

\begin{definition}
\keyterm{算术函数（arithmetic function）}是定义在正整数集上的函数，本章允许其取复数值。若 $f(1)=1$，且 $\gcd(m,n)=1$ 时 $f(mn)=f(m)f(n)$，则称 $f$ 为\keyterm{积性函数（multiplicative function）}；若乘法等式对所有正整数 $m,n$ 成立，则称完全积性。
\end{definition}

记 $\tau(n)$ 为正约数个数，$\sigma(n)$ 为正约数和；$\sum_{d\mid n}$ 总对正约数求和。由唯一分解，若 $n=\prod_{i=1}^r p_i^{\alpha_i}$，每个约数对应唯一的指数选择，故
\begin{equation}
 \tau(n)=\prod_i(\alpha_i+1),\qquad
 \sigma(n)=\prod_i(1+p_i+\cdots+p_i^{\alpha_i}).
 \label{eq:divisor-functions}
\end{equation}
$n=1$ 时均为 $1$。互素两数的素因子集合不交，所以两函数均积性；但 $\tau(4)=3\ne\tau(2)^2$，故 $\tau$ 不完全积性。一般积性函数由所有素数幂上的取值唯一确定。

以 $12=2^2\cdot3$ 为例，它的每个正约数唯一写为 $2^i3^j$，其中 $i\in\{0,1,2\}$、$j\in\{0,1\}$。共有 $3\cdot2=6$ 种选择；把约数相加则是将 $(1+2+4)(1+3)$ 展开，得到 $1+2+3+4+6+12=28$。乘积公式的根据是各个素因子指数可以独立选择，不是单纯地把一个数分成乘积后就能分开计算。

\begin{proposition}[约数求和保持积性]
\label{prop:divisor-sum-multiplicative}
若 $f$ 为积性函数，定义 $F(n)=\sum_{d\mid n}f(d)$，则 $F$ 也为积性函数。
\end{proposition}
\begin{proof}
$F(1)=f(1)=1$。若 $\gcd(m,n)=1$，则 $mn$ 的每个正约数唯一写为 $uv$，其中 $u\mid m,v\mid n$，并且 $\gcd(u,v)=1$。故
\[
 F(mn)=\sum_{u\mid m}\sum_{v\mid n}f(uv)
 =\sum_{u\mid m}\sum_{v\mid n}f(u)f(v)=F(m)F(n).
\]
\end{proof}

取 $f(d)=1$ 得到 $\tau$ 的积性，取 $f(d)=d$ 得到 $\sigma$ 的积性。一个常见误区是只知道所有 $f(p)$ 就认为确定了积性函数；实际还要知道 $f(p^2),f(p^3),\ldots$。只有完全积性函数才满足 $f(p^k)=f(p)^k$。

\section{欧拉函数与互素计数}

\begin{definition}
\keyterm{欧拉函数（Euler's totient function）}为
\[
 \varphi(n)=\#\{a:1\le a\le n,\ \gcd(a,n)=1\}.
\]
特别地，$\varphi(1)=1$。
\end{definition}

\begin{theorem}
\label{thm:totient}
欧拉函数积性，且
\begin{equation}
 \varphi(n)=n\prod_{p\mid n}\left(1-\frac1p\right),
 \qquad \sum_{d\mid n}\varphi(d)=n.
 \label{eq:totient-product}
\end{equation}
乘积遍历 $n$ 的不同素因子。
\end{theorem}
\begin{proof}
在 $1,\ldots,p^k$ 中，恰有 $p^{k-1}$ 个数被 $p$ 整除，故 $\varphi(p^k)=p^k-p^{k-1}$。若 $m,n$ 互素，中国剩余定理把模 $mn$ 的可逆类与模 $m,n$ 的可逆类对一一对应，故 $\varphi(mn)=\varphi(m)\varphi(n)$，乘积公式随之成立。

将 $1\le a\le n$ 按 $\gcd(a,n)=g$ 分类，其中 $g\mid n$。写 $a=gb$，则 $1\le b\le n/g$ 且 $\gcd(b,n/g)=1$，这一类有 $\varphi(n/g)$ 个数。各类不交且覆盖全部 $n$ 个数，得到约数求和式。
\end{proof}

\begin{example}
计算 $360$ 的约数个数、约数和与欧拉函数值。
\end{example}
\begin{solve}
$360=2^3\cdot3^2\cdot5$，所以 $\tau(360)=4\cdot3\cdot2=24$，$\sigma(360)=15\cdot13\cdot6=1170$，$\varphi(360)=360\cdot\frac12\cdot\frac23\cdot\frac45=96$。
\end{solve}

\section{狄利克雷卷积与莫比乌斯函数}

先观察 $n=6$ 的约数累计式：$F(6)=f(1)+f(2)+f(3)+f(6)$。若还知道 $F(1),F(2),F(3)$，则消去重复贡献后得到
\[
 f(6)=F(6)-F(3)-F(2)+F(1).
\]
这里加回 $F(1)$，是因为它在前两次减法中被多减了一次。这与有限集合的容斥原理相同；莫比乌斯函数将这种加减规律统一记录下来。

\begin{definition}
算术函数 $f,g$ 的\keyterm{狄利克雷卷积（Dirichlet convolution）}定义为
\begin{equation}
 (f*g)(n)=\sum_{d\mid n}f(d)g(n/d).
 \label{eq:dirichlet-convolution}
\end{equation}
记 $\mathbf1(n)=1$，$\varepsilon(1)=1$、$\varepsilon(n)=0$（$n>1$）。
\end{definition}

换元 $d\leftrightarrow n/d$ 证明交换律；两种结合方式都等于有限和 $\sum_{abc=n}f(a)g(b)h(c)$，故结合律成立；直接计算得 $f*\varepsilon=f$。若 $f,g$ 积性，互素 $m,n$ 的每个约数唯一写为 $uv$（$u\mid m,v\mid n$），将双重和分离即可证明 $f*g$ 积性。

卷积与逐点乘法不同。例如 $(f*g)(6)$ 要累加四项
\[
 f(1)g(6)+f(2)g(3)+f(3)g(2)+f(6)g(1),
\]
而逐点乘法只有 $f(6)g(6)$。恒一函数 $\mathbf1$ 的作用是进行约数累加，单位函数 $\varepsilon$ 的作用才是保持原函数不变。学习卷积时，把这两种函数区分开，比记忆形式相似的运算规则更重要。

\begin{definition}
\keyterm{莫比乌斯函数（Möbius function）}定义为
\[
 \mu(n)=
 \begin{cases}
 1,&n=1,\\
 (-1)^r,&n\text{ 是 }r\text{ 个不同素数的乘积},\\
 0,&n\text{ 被某个素数的平方整除}.
 \end{cases}
\]
\end{definition}

$\mu$ 的积性直接来自不交素因子集合。若 $n>1$ 有 $r$ 个不同素因子，只有无平方因子的约数对 $\sum_{d\mid n}\mu(d)$ 有贡献；选出 $j$ 个素因子共有 $\binom rj$ 种，故该和为 $(1-1)^r=0$。加上 $n=1$ 的情形，得到
\begin{equation}
 \mu*\mathbf1=\varepsilon.
 \label{eq:mobius-unit}
\end{equation}

\section{莫比乌斯反演及计数应用}

\begin{theorem}[莫比乌斯反演]
\label{thm:mobius-inversion}
对算术函数 $f,F$，
\begin{equation}
 F(n)=\sum_{d\mid n}f(d)
 \quad\Longleftrightarrow\quad
 f(n)=\sum_{d\mid n}\mu(d)F(n/d).
 \label{eq:mobius-inversion}
\end{equation}
\end{theorem}
\begin{proof}
左式为 $F=f*\mathbf1$，两侧与 $\mu$ 卷积，利用结合律和式\eqref{eq:mobius-unit}得 $F*\mu=f$。反之，对右式与 $\mathbf1$ 卷积即可恢复左式。所有求和均有限。
\end{proof}

\keyterm{莫比乌斯反演（Möbius inversion）}把按约数累加的信息还原。令 $F(n)=n$，由欧拉函数的约数和式得到 $\varphi(n)=\sum_{d\mid n}\mu(d)n/d$，也可再次推出乘积公式。

也可以展开一次求和，看清为什么原始贡献恰好留下。若 $F(n)=\sum_{e\mid n}f(e)$，则
\[
 \sum_{d\mid n}\mu(d)F(n/d)
 =\sum_{e\mid n}f(e)\sum_{d\mid n/e}\mu(d).
\]
内层和在 $e<n$ 时为零，在 $e=n$ 时为 $1$，所以只剩 $f(n)$。这说明反演不是神秘的形式消去，而是用权重 $\mu(d)$ 抵消所有来自真约数的重复计数。实际应用时，应先明确 $F(n)$ 是否确实按约数累加；普通的前缀和 $\sum_{k=1}^n f(k)$ 不能直接套用同一公式。

\begin{example}
求 $1\le a,b\le N$ 中互素有序数对的个数 $C(N)$，其中 $N\in\mathbb Z_{>0}$。
\end{example}
\begin{solve}
由式\eqref{eq:mobius-unit}，$\sum_{d\mid\gcd(a,b)}\mu(d)$ 在两数互素时为 $1$，否则为 $0$。交换有限求和得到
\begin{equation}
 C(N)=\sum_{d=1}^N\mu(d)\left\lfloor\frac Nd\right\rfloor^2.
 \label{eq:coprime-pairs}
\end{equation}
平方项分别选择被 $d$ 整除的两个坐标，已包括 $(1,1)$；当 $N=3$ 时结果为 $9-1-1=7$。
\end{solve}

\section{无平方因子的整数与另一种计数}

正整数若不被任何素数的平方整除，称为\keyterm{无平方因子整数（squarefree integer）}；$1$ 也属于这一类。定义表明 $\mu(n)^2$ 正好是这类整数的示性值：满足条件时为 $1$，否则为 $0$。

\begin{proposition}
\label{prop:squarefree-count}
对正整数 $n$，有
\[
 \mu(n)^2=\sum_{d^2\mid n}\mu(d).
\]
因此，不超过正整数 $N$ 的无平方因子整数个数为
\[
 \sum_{d=1}^{\lfloor\sqrt N\rfloor}\mu(d)
 \left\lfloor\frac{N}{d^2}\right\rfloor.
\]
\end{proposition}
\begin{proof}
设 $n$ 中指数至少为 $2$ 的不同素因子共有 $r$ 个。对右侧有非零贡献的 $d$，只能从这 $r$ 个素数中各取零次或一次，因此总和为 $(1-1)^r$。当 $r=0$ 时，唯一的这种 $d$ 是 $1$，总和为 $1$；当 $r>0$ 时总和为零。这正是 $\mu(n)^2$。再对 $1\le n\le N$ 求和并交换有限求和次序，每个 $d$ 对应 $\lfloor N/d^2\rfloor$ 个倍数，得到计数式。
\end{proof}

例如 $N=10$ 时，公式给出 $10-2-1=7$，对应 $1,2,3,5,6,7,10$。与互素数对计数相比，这里的共同结构是先构造取值为 $0$ 或 $1$ 的判别和，再交换求和次序。它是将整除条件转化为可计算公式的一种常用方法。

\section{习题}

积性允许沿互素因子分解计算，不能忽略互素条件。约数求和是卷积的一种，莫比乌斯反演则撤销这种累加。计数应用的关键是先构造正确的示性值，再交换有限求和；应能解释每个取整项数的是什么。

\begin{enumerate}
 \item 计算 $\varphi(840)$、$\tau(840)$ 和 $\mu(840)$。
 \item 证明 $\sum_{d\mid n}\mu(d)\tau(n/d)=1$。
 \item 证明 $\sum_{d\mid n}\mu(d)^2=2^r$，其中 $r$ 为 $n$ 的不同素因子个数。
 \item 通过计数对称性证明 $C(N)=2\sum_{k=1}^N\varphi(k)-1$，并与式\eqref{eq:coprime-pairs}比较。
 \item\label{ex:arith-noncoprime} 计算 $\varphi(12),\varphi(18),\varphi(216)$，说明为什么不能用前两者的乘积计算第三者。
 \item\label{ex:arith-invert-twelve} 若 $F(n)=\sum_{d\mid n}f(d)$，将 $f(12)$ 写成若干个 $F$ 值的加减，并逐项解释系数。
 \item\label{ex:arith-squarefree} 用命题\ref{prop:squarefree-count}计算不超过 $20$ 的无平方因子整数个数，再用枚举核对。
 \item\label{ex:arith-inverse} 证明算术函数 $f$ 存在卷积逆元 $g$（即 $f*g=\varepsilon$）当且仅当 $f(1)\ne0$，且逆元唯一。给出按 $n$ 递增求 $g(n)$ 的递推式。
\end{enumerate}

\chapter{费马、欧拉与威尔逊定理}
\label{ch:classical-congruences}

计算 $2^{1000}$ 的完整十进制展开很费力，但求它除以 $7$ 的余数只需抓住幂的周期。本章的费马、欧拉定理给出保证出现的周期，威尔逊定理则用阶乘刻画素数。三个定理都来自有限剩余类的结构，不过适用条件和逻辑方向不同；理解它们的证明，能够减少套公式时最常见的错误。

\section{费马小定理及两种表述}

\begin{theorem}[费马小定理]
\label{thm:fermat}
若 $p$ 为素数且 $p\nmid a$，则 $a^{p-1}\equiv1\pmod p$。等价地，对任意整数 $a$，$a^p\equiv a\pmod p$。
\end{theorem}
\begin{proof}
$a,2a,\ldots,(p-1)a$ 均非零，且模 $p$ 两两不同，因为 $ia\equiv ja$ 可约去 $a$。故它们是 $1,\ldots,p-1$ 的一个排列，从而
\[
 a^{p-1}(p-1)!\equiv(p-1)!\pmod p.
\]
阶乘与 $p$ 互素，约去它便得结论。再乘以 $a$ 得第二种形式；$p\mid a$ 时第二式两侧均为零。反之，在 $p\nmid a$ 时第二式可约去 $a$。
\end{proof}

\keyterm{费马小定理（Fermat's little theorem）}控制模素数的幂。例如 $2^6\equiv1\pmod7$，而 $1000\equiv4\pmod6$，故 $2^{1000}\equiv2^4\equiv2\pmod7$。

证明中的“排列”可以在模 $7$、$a=3$ 时直接看见：$3,6,9,12,15,18$ 的最小正余数依次为 $3,6,2,5,1,4$，只是把 $1,2,3,4,5,6$ 换了次序。相乘会多出六个因子 $3$，但余数的总乘积不变，所以可约去 $6!$ 得到 $3^6\equiv1\pmod7$。这是利用可逆乘法重新排列有限集合的证明模式，欧拉定理将把它推广到任意模数的可逆类。

\section{欧拉定理与指数化简}

\begin{theorem}[欧拉定理]
\label{thm:euler}
若 $m\ge2$ 且 $\gcd(a,m)=1$，则 $a^{\varphi(m)}\equiv1\pmod m$。
\end{theorem}
\begin{proof}
取一个简化剩余系 $r_1,\ldots,r_{\varphi(m)}$。乘以 $a$ 后各项仍可逆，且由约分法则两两不同，故形成相同类的排列。两侧乘积同余，并可约去与 $m$ 互素的 $\prod_i r_i$，得到结论。
\end{proof}

\keyterm{欧拉定理（Euler's theorem）}包含费马小定理。若非负指数 $u\equiv v\pmod{\varphi(m)}$ 且底数与 $m$ 互素，则 $a^u\equiv a^v\pmod m$。去掉互素条件则不成立：$\varphi(8)=4$，但 $2^4\equiv0\not\equiv1=2^0\pmod8$。

处理不互素底数时，可先按素数幂拆分模数；对整除底数的素因子，用指数增长判断何时变为零，对其他素数幂再应用欧拉定理，最后用中国剩余定理合并。

\begin{example}[底数与模数不互素时的幂运算]
求 $2^{100}$ 除以 $100$ 的余数。
\end{example}
\begin{solve}
因为 $\gcd(2,100)\ne1$，不能直接按 $\varphi(100)=40$ 对指数约化。拆成 $100=4\cdot25$：模 $4$ 时，指数至少为 $2$，故 $2^{100}\equiv0$；模 $25$ 时，$\gcd(2,25)=1$ 且 $\varphi(25)=20$，故 $2^{100}\equiv1$。写所求余数 $r=1+25t$，模 $4$ 得 $1+t\equiv0$，所以 $t\equiv3\pmod4$，最终 $r\equiv76\pmod{100}$。最小非负余数为 $76$。
\end{solve}

\begin{example}[多层幂的模数需要逐层确定]
求 $3^{,3^{10}}\pmod{17}$。
\end{example}
\begin{solve}
外层底数 $3$ 与 $17$ 互素，所以只需先求指数 $3^{10}$ 模 $16$ 的值。由 $3^4\equiv1\pmod{16}$ 得 $3^{10}\equiv9\pmod{16}$；因此原式同余于 $3^9$。模 $17$ 依次平方得 $3^2\equiv9$、$3^4\equiv13$、$3^8\equiv16$，于是 $3^9\equiv14$。这里内层指数所用的模数是外层的周期 $16$，而不是外层模数 $17$。
\end{solve}

\section{威尔逊定理与逆命题辨析}

\begin{theorem}[威尔逊定理]
\label{thm:wilson}
整数 $n>1$ 为素数，当且仅当 $(n-1)!\equiv-1\pmod n$。
\end{theorem}
\begin{proof}
若 $n=p$ 为奇素数，在非零剩余类中，将每个元素与其逆元配对。满足自逆的元素恰为 $1,-1$：$x^2\equiv1$ 意味着 $p\mid(x-1)(x+1)$，由欧几里得引理得 $x\equiv\pm1$。其他配对乘积均为 $1$，所以总乘积为 $-1$。$p=2$ 时直接验证。

反之，若 $n$ 合数，取 $1<d<n$ 且 $d\mid n$。因为 $d$ 是阶乘中的一项，$d\mid(n-1)!$；所给同余又使 $d\mid((n-1)!+1)$，于是 $d\mid1$，矛盾。
\end{proof}

\keyterm{威尔逊定理（Wilson's theorem）}是素性的充要条件，但直接计算长阶乘并不高效。费马同余只是必要条件：$341=11\cdot31$ 为合数，而 $2^{10}=1024\equiv1\pmod{341}$，故 $2^{340}\equiv1\pmod{341}$。

例如模 $7$ 的逆元配对为 $(2,4)$、$(3,5)$，两对乘积都余 $1$，剩下自逆的 $1$ 与 $6=-1$，故 $6!\equiv-1$。逆命题则不需要逐个分析所有合数的阶乘余数，只要指出一个真因子会同时整除阶乘和阶乘加一，就已产生矛盾。数学上存在充要判据，与计算上能够迅速应用判据，是不同的问题。

\section{快速模幂与计算验证}

\keyterm{快速模幂（modular exponentiation by repeated squaring）}直接计算 $a^e\pmod m$，适用于任意整数 $a$、非负整数 $e$ 和 $m\ge2$。初始化 $R=1$、$B$ 为 $a$ 的最小非负余数、$E=e$；在 $E>0$ 时重复：
\begin{enumerate}
 \item 若 $E$ 为奇数，令 $R$ 为 $RB$ 模 $m$ 的最小非负余数。
 \item 令 $B$ 为 $B^2$ 模 $m$ 的最小非负余数，令 $E\leftarrow\lfloor E/2\rfloor$。
\end{enumerate}
循环不变式为 $RB^E\equiv a^e\pmod m$。偶数步将 $B^E$ 写成 $(B^2)^{E/2}$；奇数步先提取一个 $B$，因此不变式保持。$E$ 每步约减半，最终为零，此时 $R$ 即所求。$e=0$ 时直接返回 $1$。

例如计算 $3^{13}\pmod{17}$，平方得到 $3^2\equiv9$、$3^4\equiv13$、$3^8\equiv16$；由 $13=8+4+1$ 得 $3^{13}\equiv16\cdot13\cdot3\equiv12\pmod{17}$。模幂算法本身不需要互素条件。

\section{习题}

费马与欧拉定理给出可逆底数的幂周期，威尔逊定理刻画素性，快速模幂则直接执行计算。约化指数前先确认底数可逆；若不互素，可拆分模数或直接采用快速模幂，不应机械套用周期。

\begin{enumerate}
 \item 求 $7^{222}\pmod{13}$ 与 $3^{100}\pmod{20}$。
 \item 证明对素数 $p>3$，$p^2\equiv1\pmod{24}$。
 \item 若奇素数 $p\equiv1\pmod4$，证明 $\bigl(((p-1)/2)!\bigr)^2\equiv-1\pmod p$。
 \item 用快速模幂计算 $5^{117}\pmod{19}$，并用欧拉定理核对。
 \item\label{ex:classical-noncoprime} 求 $6^{100}\pmod{35}$ 与 $6^{100}\pmod{40}$，分别说明所用周期或局部整除条件。
 \item\label{ex:classical-factorial} 求 $10!\pmod{13}$。提示：把缺少的两项补成 $12!$。
 \item\label{ex:classical-fermat-all} 用费马小定理证明每个整数 $a$ 都满足 $30\mid a^5-a$，注意分别处理底数被素数整除的情形。
 \item\label{ex:classical-pseudoprime} 验证 $341$ 通过底数 $2$ 的费马检验，再用模 $31$ 的计算说明它不能通过底数 $3$ 的费马检验。
\end{enumerate}

\chapter{阶、原根与高次同余}
\label{ch:primitive-roots}

欧拉定理保证 $a^{\varphi(m)}\equiv1\pmod m$，但第一次回到 $1$ 往往更早。例如 $2^3\equiv1\pmod7$，而 $\varphi(7)=6$。最短周期称为乘法阶；若这个周期恰好遍历所有可逆类，对应的底数就是原根。原根的价值在于给乘法建立一个指数坐标，从而把高次同余转化为一次同余。

\section{乘法阶及其整除性质}

\begin{definition}
若 $m\ge2$ 且 $\gcd(a,m)=1$，使 $a^t\equiv1\pmod m$ 成立的最小正整数 $t$ 称为 $a$ 模 $m$ 的\keyterm{乘法阶（multiplicative order）}，记为 $\operatorname{ord}_m(a)$。
\end{definition}

欧拉定理保证至少有一个这样的正指数，因此最小值存在。

模 $7$ 时，$2^0,2^1,2^2,2^3,\ldots$ 的余数为 $1,2,4,1,\ldots$，故 $\operatorname{ord}_7(2)=3$；底数 $3$ 的幂却要经过 $1,3,2,6,4,5$ 六个类后才回到 $1$。阶依赖底数和模数两者，不能只从 $m$ 推断所有底数的阶都等于 $\varphi(m)$。此外，若 $a$ 不可逆，就不存在 $a^t\equiv1$ 的正指数，否则 $a^{t-1}$ 已给出逆元，所以定义中的互素条件不可省略。

\begin{proposition}
\label{prop:order}
令 $t=\operatorname{ord}_m(a)$。对正整数 $k$，$a^k\equiv1\pmod m$ 当且仅当 $t\mid k$；特别地，$t\mid\varphi(m)$。此外
\begin{equation}
 \operatorname{ord}_m(a^k)=\frac{t}{\gcd(t,k)}.
 \label{eq:order-power}
\end{equation}
\end{proposition}
\begin{proof}
写 $k=qt+r$、$0\le r<t$，则 $a^k\equiv a^r$；由 $t$ 的最小性，$a^k\equiv1$ 当且仅当 $r=0$。对第二式，$(a^k)^s\equiv1$ 等价于 $t\mid ks$，满足该条件的最小正整数 $s$ 为 $t/\gcd(t,k)$。
\end{proof}

若 $b>1$ 且 $\gcd(b,10)=1$，长除法产生的余数序列是 $10^j\pmod b$；首次回到余数 $1$ 的时间恰为 $\operatorname{ord}_b(10)$，这就是 $1/b$ 的十进制最小循环节长度。例如 $10\equiv3\pmod7$，$3$ 模 $7$ 的阶为 $6$，故 $1/7$ 的循环节长为 $6$。

\begin{example}[先列周期的约数，再排除]
求 $\operatorname{ord}_{17}(4)$。
\end{example}
\begin{solve}
阶必须整除 $16$，候选值只有 $1,2,4,8,16$。计算 $4^2\equiv-1\pmod{17}$，得 $4^4\equiv1$，而 $4\not\equiv1$、$4^2\not\equiv1$，所以阶为 $4$。先用整除关系缩小范围，比逐项计算直至重复更有条理。
\end{solve}

更一般地，对与 $10$ 互素的 $b>1$ 及 $1\le a<b$、$\gcd(a,b)=1$，分数 $a/b$ 的最小循环节也为 $\operatorname{ord}_b(10)$。长除法每一步的余数满足 $r_{j+1}\equiv10r_j\pmod b$，初值 $r_0=a$；回到初始余数等价于 $10^ja\equiv a$，约去 $a$ 后便得到阶的条件。这说明同分母的既约分数可以有不同的循环数位，却有相同的循环长度。

\section{原根的含义与存在范围}

\begin{definition}
若 $\operatorname{ord}_m(g)=\varphi(m)$，称 $g$ 为模 $m$ 的\keyterm{原根（primitive root）}。
\end{definition}

此时 $g^0,g^1,\ldots,g^{\varphi(m)-1}$ 两两不同，恰好遍历全部可逆类。模 $7$ 时，$3$ 的幂依次为 $1,3,2,6,4,5$，故 $3$ 是原根；模 $8$ 的每个奇数平方均为 $1$，所以没有阶为 $\varphi(8)=4$ 的元素。

\begin{proposition}[原根的实用检验]
\label{prop:primitive-root-test}
设 $\gcd(g,m)=1$，$h=\varphi(m)>1$。则 $g$ 为模 $m$ 的原根，当且仅当对 $h$ 的每个不同素因子 $q$，都有 $g^{h/q}\not\equiv1\pmod m$。
\end{proposition}
\begin{proof}
若阶为 $h$，任何正整数 $h/q<h$ 都不能使幂为 $1$。反之，设阶 $t$ 是 $h$ 的真约数，取整数 $h/t>1$ 的一个素因子 $q$，则 $t\mid h/q$，所以 $g^{h/q}\equiv1$，与条件矛盾。
\end{proof}

例如检验 $2$ 是否为模 $11$ 的原根，只需检查 $10$ 的不同素因子 $2,5$：$2^5\equiv-1$、$2^2\equiv4$，均不为 $1$，所以阶为 $10$。这个方法需要先知道 $\varphi(m)$ 的素因子分解；它是候选原根的检验方法，本身并不证明每个模数都有原根。

\begin{lemma}[模素数的根数界]
\label{lem:polynomial-roots}
若整系数多项式 $f$ 模素数 $p$ 后为次数 $d$ 的非零多项式，则 $f(x)\equiv0\pmod p$ 至多有 $d$ 个根类。
\end{lemma}
\begin{proof}
对 $d$ 归纳。非零常数无根。若 $a$ 是根，注意每个差 $x^j-a^j$ 都被多项式 $x-a$ 整除，因为
\[
 x^j-a^j=(x-a)(x^{j-1}+ax^{j-2}+\cdots+a^{j-1}).
\]
因此 $f(x)-f(a)=(x-a)g(x)$；将系数模 $p$ 约化并使用 $f(a)\equiv0$，得到 $f(x)=(x-a)g(x)$，其中 $g$ 次数为 $d-1$。任意不同于 $a$ 的根 $b$ 都满足 $(b-a)g(b)\equiv0$，由于 $b-a$ 可逆，$g(b)\equiv0$。故除 $a$ 外至多还有 $d-1$ 个根，归纳完成。
\end{proof}

\begin{theorem}
\label{thm:prime-primitive-root}
每个素数模数都有原根。
\end{theorem}
\begin{proof}
目标是构造阶恰为 $p-1$ 的元素。直接寻找可能不容易，但可以先对 $p-1$ 中的每个素数幂构造对应阶的元素，再利用互素性把这些局部周期相乘。

$p=2$ 时取 $1$。对奇素数 $p$，写 $p-1=\prod_i\ell_i^{e_i}$。对每个 $i$，多项式 $X^{(p-1)/\ell_i}-1$ 的根数小于 $p-1$，故存在非零类 $a_i$ 不满足该方程。由阶整除 $p-1$，$\operatorname{ord}_p(a_i)$ 的 $\ell_i$ 因子指数必为 $e_i$。令 $b_i=a_i^{(p-1)/\ell_i^{e_i}}$，由式\eqref{eq:order-power}，$b_i$ 的阶为 $\ell_i^{e_i}$。

若两个可逆类 $u,v$ 的阶为互素的 $r,s$，则 $uv$ 的阶为 $rs$。事实上，$(uv)^k=1$ 给出 $u^k=v^{-k}$，这一公共类的阶同时整除 $r,s$，故它为 $1$，于是 $r\mid k,s\mid k$；反向 $(uv)^{rs}=1$。逐次应用可知 $\prod_i b_i$ 的阶为 $p-1$。
\end{proof}

\begin{theorem}[原根分类]
\label{thm:primitive-root-classification}
整数 $m\ge2$ 存在原根，当且仅当 $m=2,4,p^k,2p^k$，其中 $p$ 为奇素数、$k\ge1$。
\end{theorem}
\begin{proof}
证明分为“这些模数确实有原根”和“其余模数没有原根”。前一部分把模 $p$ 的原根提升到 $p^k$；后一部分把模数拆开，证明不同部分的偶周期会重叠，使总周期达不到全部可逆类的个数。

先证明存在性。$2,4$ 分别取 $1,3$。若 $g$ 为模奇素数 $p$ 的原根，则 $g$ 与 $g+p$ 中至少一个数 $h$ 满足 $h^{p-1}\not\equiv1\pmod{p^2}$，因为二者的 $(p-1)$ 次幂之差模 $p^2$ 为 $(p-1)g^{p-2}p\not\equiv0$。故 $v_p(h^{p-1}-1)=1$。

对 $p\nmid c$、$s\ge1$，二项展开给出
\[
 (1+cp^s)^p\equiv1+cp^{s+1}\pmod{p^{s+2}}.
\]
这里使用 $p$ 为奇素数。归纳得到 $v_p(h^{(p-1)p^j}-1)=j+1$。模 $p^k$ 的阶既是 $p-1$ 的倍数，又整除 $(p-1)p^{k-1}$；上述指数关系迫使它等于 $(p-1)p^{k-1}$，故 $h$ 为原根。取同一模 $p^k$ 类的奇数代表元，由中国剩余定理，其模 $2p^k$ 的阶不变。

具体地，模 $p^k$ 的阶必形如 $(p-1)p^r$，其中 $0\le r\le k-1$。若 $r<k-1$，上式给出 $v_p(h^{(p-1)p^r}-1)=r+1<k$，这与该幂模 $p^k$ 为 $1$ 矛盾。这一步说明提升后的周期必须达到最大值，而不只是“比原来更长”。

再证必要性。若 $m=uv$，其中 $u,v>2$ 互素，则可逆类的阶整除 $\operatorname{lcm}(\varphi(u),\varphi(v))$；两欧拉函数值均为偶数，所以该最小公倍数严格小于 $\varphi(u)\varphi(v)=\varphi(m)$。这里 $\varphi(n)$ 在 $n>2$ 时为偶数，是因为可逆类可不重复地配成 $a,-a$。于是 $m$ 不能同时含两个不同的奇素因子，也不能既被 $4$ 整除又有奇素因子。余下只需排除 $m=2^k$、$k\ge3$：奇数 $a$ 满足 $a^2\equiv1\pmod8$，逐次平方得 $a^{2^{k-2}}\equiv1\pmod{2^k}$，所以阶小于 $\varphi(2^k)$。
\end{proof}

\section{指标与幂同余方程}

固定模 $m$ 的原根 $g$，令 $h=\varphi(m)$。每个可逆类 $a$ 唯一写为 $a\equiv g^s\pmod m$，其中 $s$ 在模 $h$ 意义下唯一，称为关于 $g$ 的\keyterm{指标（index）}或\keyterm{离散对数（discrete logarithm）}。指标把可逆类的乘法转化为模 $h$ 的加法。

\begin{proposition}
若模 $m$ 有原根，$k\ge1$、$\gcd(a,m)=1$，令 $d=\gcd(k,\varphi(m))$，则 $x^k\equiv a\pmod m$ 可解当且仅当 $a^{\varphi(m)/d}\equiv1\pmod m$；可解时恰有 $d$ 个根类。
\end{proposition}
\begin{proof}
任何解 $x$ 必可逆，否则其与 $m$ 的某个公共素因子也整除 $a$。写 $x=g^t,a=g^s$，原方程等价于 $kt\equiv s\pmod h$，可解当且仅当 $d\mid s$，且恰有 $d$ 个指数类。另一方面，$a^{h/d}=1$ 等价于 $h\mid sh/d$，也就是 $d\mid s$。
\end{proof}

例如模 $7$ 取原根 $3$，$2\equiv3^2$，则 $x^2\equiv2$ 等价于 $2t\equiv2\pmod6$，得到 $t\equiv1,4\pmod6$，即 $x\equiv3,4\pmod7$。无原根的模数不能直接使用这一化简。

\begin{example}[先求指数，再还原根]
求 $x^3\equiv8\pmod{13}$ 的全部根类。
\end{example}
\begin{solve}
$2^6\equiv-1$、$2^4\equiv3\pmod{13}$，由命题\ref{prop:primitive-root-test}，$2$ 是模 $13$ 的原根。因 $8=2^3$，写 $x\equiv2^t$，原式等价于 $3t\equiv3\pmod{12}$，故 $t\equiv1,5,9\pmod{12}$。还原得到 $x\equiv2,6,5\pmod{13}$。三个指数类不同，原根的幂也不同，所以恰有三根；只找出 $x=2$ 会遗漏另外两根。
\end{solve}

\section{多项式同余与解的提升（选学）}

已知模 $p^j$ 的一个根后，可以尝试保留已有余数，再寻找下一位修正。下述命题中的 $f'(x)$ 是多项式的形式导数：若 $f(x)=\sum_i c_ix^i$，则 $f'(x)=\sum_{i\ge1}ic_ix^{i-1}$。所需恒等式由二项展开即可推出，不要求先学习微积分。

\begin{proposition}[简单根的亨泽尔提升]
\label{prop:hensel}
设 $f\in\mathbb Z[x]$，$p$ 为素数，$j\ge1$，$f(r)\equiv0\pmod{p^j}$ 且 $f'(r)\not\equiv0\pmod p$。则根类 $r\pmod{p^j}$ 唯一提升为模 $p^{j+1}$ 的根类。
\end{proposition}
\begin{proof}
所有提升写为 $r+tp^j$，其中 $t\pmod p$。由二项展开，
\[
 f(r+tp^j)\equiv f(r)+tp^jf'(r)\pmod{p^{j+1}},
\]
因为高次项含 $p^{2j}$，而 $2j\ge j+1$。故所求条件为 $f(r)/p^j+tf'(r)\equiv0\pmod p$，其系数可逆，所以 $t$ 唯一。
\end{proof}

这称为\keyterm{亨泽尔提升（Hensel lifting）}。例如 $f(x)=x^2-2$ 的模 $7$ 根 $3$ 提升至模 $49$ 时，$f(3)/7=1$、$f'(3)=6$，故 $1+6t\equiv0\pmod7$，得到 $t=1$，根为 $10$。若导数为零，则可能无提升，也可能有多个提升，需重新求解。

将 $r$ 替换为 $r+tp^j$，等于保留已有的模 $p^j$ 信息，只寻找下一位的修正量 $t$。例如模 $3$ 时 $x=0$ 同时是 $x^2\equiv0$ 与 $x^2\equiv3$ 的根；提升到模 $9$，前者有 $0,3,6$ 三个提升，后者一个也没有。这正说明导数不为零的条件保证的是唯一提升。

\section{习题}

乘法阶是最短正周期，原根使幂恰好遍历全部可逆类。求阶时先检查欧拉函数的约数，解幂同余时先确认原根存在，再将根与指数一一对应。提升根的唯一性还需单独检查导数条件。

\begin{enumerate}
 \item 求模 $7$ 的全部原根，证明存在原根的模数 $m$ 恰有 $\varphi(\varphi(m))$ 个原根类。
 \item 求 $1/13$ 的最小循环节长度。
 \item 求 $x^3\equiv1\pmod7$ 的全部解。
 \item 将 $x^2\equiv2\pmod7$ 的两个根提升至模 $49$，再代入核查。
 \item\label{ex:order-list} 求模 $9$ 的每个可逆类的阶，并列出全部原根。
 \item\label{ex:order-composite} 设整数 $m,n\ge2$。若 $\gcd(m,n)=1$ 且 $\gcd(a,mn)=1$，证明 $\operatorname{ord}_{mn}(a)=\operatorname{lcm}(\operatorname{ord}_m(a),\operatorname{ord}_n(a))$。
 \item\label{ex:order-sixth} 求 $x^6\equiv1\pmod{13}$ 的全部根类。
 \item\label{ex:order-singular} 比较模 $3$ 的根 $0$ 在 $x^2\equiv0$ 与 $x^2\equiv3$ 中提升到模 $9$ 的结果；再求 $x^2\equiv0\pmod{27}$ 的全部根类。
\end{enumerate}

\chapter{二次剩余与二次互反律}
\label{ch:quadratic-residues}

一次同余可以用逆元求解，二次同余却增加了一个新的问题：并非每个非零类都是平方。例如模 $7$ 的非零平方只有 $1,2,4$，因此 $x^2\equiv3\pmod7$ 无解。本章先建立“是否为平方”的判别符号，再用二次互反律把大模数下的判别归结到较小模数。应始终区分判定可解性、计算平方根、计算根的个数这三个任务。

\section{二次剩余与勒让德符号}

设 $m\ge2$。本章在与 $m$ 互素的整数中讨论\keyterm{二次剩余（quadratic residue）}：若 $\gcd(a,m)=1$ 且 $x^2\equiv a\pmod m$ 有解，称 $a$ 为二次剩余；若 $\gcd(a,m)=1$ 而该方程无解，称二次非剩余。与 $m$ 不互素的情形另行讨论，不把它们一概称为二次非剩余。以下先固定奇素数 $p$。

\begin{definition}
\keyterm{勒让德符号（Legendre symbol）}定义为
\[
 \left(\frac ap\right)=
 \begin{cases}
 0,&p\mid a,\\
 1,&a\text{ 是模 }p\text{ 的二次剩余},\\
 -1,&a\text{ 是模 }p\text{ 的二次非剩余}.
 \end{cases}
\]
\end{definition}

若 $x^2\equiv y^2\pmod p$，则 $p\mid(x-y)(x+y)$，所以 $x\equiv\pm y$。非零的 $p-1$ 个类可分为 $(p-1)/2$ 对 $\{x,-x\}$，每对给出不同的平方。因此非零平方剩余恰有 $(p-1)/2$ 个，各有两个平方根类。

以 $p=11$ 为例，只需计算 $1,2,3,4,5$ 的平方，便得到全部非零平方剩余 $1,4,9,5,3$；$6,7,8,9,10$ 的平方会重复同一列表。于是 $\left(\frac5{11}\right)=1$、$\left(\frac2{11}\right)=-1$、$\left(\frac{22}{11}\right)=0$。勒让德符号看起来像分数，但其值只可能是 $-1,0,1$，不执行分数的约分运算；分母必须是奇素数。

\section{欧拉判别法与高斯引理}

\begin{theorem}[欧拉判别法]
\label{thm:euler-criterion}
对奇素数 $p$ 与整数 $a$，
\begin{equation}
 a^{(p-1)/2}\equiv\left(\frac ap\right)\pmod p.
 \label{eq:euler-criterion}
\end{equation}
\end{theorem}
\begin{proof}
$p\mid a$ 时成立。若 $a=b^2\not\equiv0$，由费马小定理得 $a^{(p-1)/2}=b^{p-1}\equiv1$。多项式 $X^{(p-1)/2}-1$ 至多有 $(p-1)/2$ 个根，已由平方剩余占满，故非剩余不能使此值为 $1$。又由费马小定理，此值的平方为 $1$，所以只能为 $-1$。
\end{proof}

\keyterm{欧拉判别法（Euler's criterion）}立即给出乘法性
\[
 \left(\frac{ab}{p}\right)=\left(\frac ap\right)\left(\frac bp\right).
\]
先由幂运算得到两侧模 $p$ 同余，再利用它们只取 $0,1,-1$ 且 $p$ 为奇素数，得到整数意义下的相等。

乘法性表明，两个非剩余的乘积是剩余；剩余与非剩余的乘积是非剩余。但勒让德符号没有相应的加法性。例如模 $7$ 时 $1$ 和 $2$ 都是平方剩余，$1+2=3$ 却不是。实际计算应先对分子作乘法分解，再分别处理各因子，不能把分子上的加法直接移到符号外。

\begin{lemma}[高斯引理（Gauss's lemma）]
\label{lem:gauss}
若 $p\nmid a$，在 $a,2a,\ldots,\frac{p-1}{2}a$ 的最小正余数中，大于 $p/2$ 的共有 $N$ 个，则 $\left(\frac ap\right)=(-1)^N$。
\end{lemma}
\begin{proof}
记 $h=(p-1)/2$，将每个余数 $r>p/2$ 替换为 $r-p$，得到绝对值均在 $1,\ldots,h$ 中的带符号余数。若两个绝对值相同，则 $ia\equiv\pm ja$，进而 $i\equiv\pm j\pmod p$；由 $1\le i,j\le h$ 只能有 $i=j$。故这些绝对值恰为 $1,\ldots,h$ 的排列。相乘并约去 $h!$，得 $a^h\equiv(-1)^N$，再用欧拉判别法。
\end{proof}

\section{二次互反律及补充公式}

在进入互反律之前，先把高斯引理落实到一次小计算。对 $p=11,a=3$，$3,6,9,12,15$ 的最小正余数为 $3,6,9,1,4$；超过 $11/2$ 的是 $6,9$，共两项，所以 $\left(\frac3{11}\right)=1$。将这两项改写为 $-5,-2$ 后，五个带符号余数的绝对值恰为 $1,2,3,4,5$。证明中“只差一个符号”的现象，在这里就是负号的数目决定最终结果。

二次互反律把 $p$ 是否为模 $q$ 的平方，与 $q$ 是否为模 $p$ 的平方联系起来。它并不是说两者总相同，而是说差别完全由 $p,q$ 模 $4$ 的余数决定。下面的证明先把符号写成格点计数的奇偶性，再利用同一个矩形被一条直线分成两部分这一事实。

\begin{theorem}[二次互反律]
\label{thm:quadratic-reciprocity}
若 $p,q$ 为不同奇素数，则
\begin{equation}
 \left(\frac pq\right)\left(\frac qp\right)
 =(-1)^{\frac{p-1}{2}\frac{q-1}{2}}.
 \label{eq:quadratic-reciprocity}
\end{equation}
对任意奇素数 $p$，
\begin{equation}
 \left(\frac{-1}{p}\right)=(-1)^{(p-1)/2},\qquad
 \left(\frac2p\right)=(-1)^{(p^2-1)/8}.
 \label{eq:quadratic-supplements}
\end{equation}
\end{theorem}
\begin{proof}
先把高斯引理改写为计数形式。令 $h=(p-1)/2$，写 $qj=pk_j+r_j$（$1\le r_j<p$），并记 $s_j=\min(r_j,p-r_j)$。由高斯引理的证明，$s_j$ 是 $1,\ldots,h$ 的排列。若 $N$ 个 $r_j>p/2$，则
\[
 \sum_jr_j=\sum_js_j+Np-2\sum_{r_j>p/2}s_j.
\]
因 $p,q$ 都为奇数，将 $q\sum_jj=p\sum_jk_j+\sum_jr_j$ 模 $2$ 化简，得到 $N\equiv\sum_jk_j\pmod2$。所以
\[
 \left(\frac qp\right)=(-1)^{\sum_{j=1}^{(p-1)/2}\lfloor qj/p\rfloor}.
\]

考虑矩形中的整数点 $1\le j\le(p-1)/2$、$1\le k\le(q-1)/2$。直线 $qj=pk$ 不经过任何这样的点，因为 $p,q$ 是不同素数。直线下方的点数为 $\sum_j\lfloor qj/p\rfloor$，上方的点数为 $\sum_k\lfloor pk/q\rfloor$。两者之和为 $(p-1)(q-1)/4$，这给出式\eqref{eq:quadratic-reciprocity}。

第一条补充公式由欧拉判别法取 $a=-1$ 得到。第二条用高斯引理取 $a=2$：余数是 $2,4,\ldots,p-1$，故 $N=(p-1)/2-\lfloor p/4\rfloor$。按 $p\equiv1,3,5,7\pmod8$ 分别计算，其奇偶性依次为偶、奇、奇、偶，与 $(p^2-1)/8$ 相同。
\end{proof}

\keyterm{二次互反律（quadratic reciprocity law）}将大分母的符号化为较小分母的符号。两个奇素数不同时模 $4$ 余 $3$ 时，交换分子分母不变号；二者都余 $3$ 时变号。

\begin{example}
判断 $x^2\equiv5\pmod{11}$ 和 $x^2\equiv3\pmod{11}$ 是否可解。
\end{example}
\begin{solve}
$\left(\frac5{11}\right)=\left(\frac{11}5\right)=\left(\frac15\right)=1$，且 $4^2\equiv5$，故根为 $4,7$。又 $\left(\frac3{11}\right)=-\left(\frac{11}3\right)=-\left(\frac23\right)=1$，根为 $5,6$。判别可解性不等于已求出平方根，还需额外计算。
\end{solve}

\section{雅可比符号与合数模数}

\begin{example}[分解、互反与取余配合使用]
计算 $\left(\frac{10}{13}\right)$ 与 $\left(\frac3{19}\right)$，据此判断相应二次同余是否可解。
\end{example}
\begin{solve}
由乘法性及补充公式，$\left(\frac2{13}\right)=-1$；又因 $5\equiv1\pmod4$，有 $\left(\frac5{13}\right)=\left(\frac{13}5\right)=\left(\frac35\right)=-1$。故 $\left(\frac{10}{13}\right)=1$，而直接计算 $6^2\equiv10\pmod{13}$ 得到两根 $6,7$。对第二个符号，$3,19$ 都模 $4$ 余 $3$，所以
\[
 \left(\frac3{19}\right)=-\left(\frac{19}3\right)
 =-\left(\frac13\right)=-1,
\]
故 $x^2\equiv3\pmod{19}$ 无解。符号变号的条件必须在每次互换时重新检查。
\end{solve}

若奇正整数 $n>1$ 分解为 $n=\prod_i p_i^{e_i}$，定义\keyterm{雅可比符号（Jacobi symbol）}
\[
 \left(\frac an\right)=\prod_i\left(\frac a{p_i}\right)^{e_i},
 \qquad \left(\frac a1\right)=1.
\]
分母为素数时与勒让德符号一致。它对分子和奇正分母均具乘法性；对互素奇正整数 $m,n$，满足同样的互反式 $\left(\frac mn\right)\left(\frac nm\right)=(-1)^{(m-1)(n-1)/4}$。这是将素数互反律逐对相乘的结果：对奇数 $u,v$，$(uv-1)/2\equiv(u-1)/2+(v-1)/2\pmod2$，故各素因子的符号指数合并为所示指数。

如果 $a$ 与 $n$ 互素且模 $n$ 为平方，则雅可比符号为 $1$；反向不成立。例如 $\left(\frac2{15}\right)=\left(\frac23\right)\left(\frac25\right)=(-1)(-1)=1$，但模 $3$ 已无平方根。模奇素数的非零平方根可用命题\ref{prop:hensel}逐次唯一提升至该素数的幂，合数模数则继续用中国剩余定理合并；模 $2$ 的幂须单独处理。

\begin{proposition}[奇素数幂上的可逆平方]
\label{prop:odd-prime-square-roots}
若 $p$ 为奇素数、$k\ge1$ 且 $p\nmid a$，则 $x^2\equiv a\pmod{p^k}$ 可解当且仅当 $\left(\frac ap\right)=1$；可解时恰有两个根类。
\end{proposition}
\begin{proof}
模 $p^k$ 的解约化后必是模 $p$ 的解，所以条件必要。反之，模 $p$ 恰有两个非零根。对 $f(x)=x^2-a$，每个根 $r$ 满足 $f'(r)=2r\not\equiv0\pmod p$，故由命题\ref{prop:hensel}逐次唯一提升到模 $p^k$。任何高模数根都必须约化为原来两根之一，而每次提升唯一，所以没有第三个根。
\end{proof}

对奇数 $n=\prod_{i=1}^r p_i^{e_i}$ 及 $\gcd(a,n)=1$，因此必须逐个检查所有 $\left(\frac a{p_i}\right)$，而不能只检查它们的带重数乘积。若全部为 $1$，中国剩余定理给出恰好 $2^r$ 个根类；只要一个为 $-1$，就无解。

\section{模二的幂的平方与一般平方同余}

奇素数的证明用到了 $2$ 可逆，因而不能直接搬到模 $2^k$ 的情形。最初几个模数给出正确的提示：奇数平方模 $4$ 恒为 $1$，模 $8$ 也恒为 $1$，而 $x^2\equiv1\pmod8$ 有四个奇数根。

\begin{proposition}
\label{prop:two-power-squares}
设 $a$ 为奇数。模 $2$ 时 $x^2\equiv a$ 恰有一个根类；模 $4$ 时可解当且仅当 $a\equiv1\pmod4$，可解时恰有两个根类。对 $k\ge3$，$x^2\equiv a\pmod{2^k}$ 可解当且仅当 $a\equiv1\pmod8$，可解时恰有四个根类。
\end{proposition}
\begin{proof}
前两种情形逐个列出余数即可。若 $k\ge3$ 且有根，根必为奇数，因此平方模 $8$ 为 $1$，必要性成立。为证充分性和根数，对 $k$ 归纳：$k=3$ 时四个奇数类都是根。

设模 $2^k$ 已有四根。将根 $r$ 与 $r+2^{k-1}$ 配对；后者仍是模 $2^k$ 的根，且与前者不同。二者平方之差为
\[
 (r+2^{k-1})^2-r^2=2^kr+2^{2k-2}\equiv2^k\pmod{2^{k+1}},
\]
这里使用 $r$ 为奇数及 $k\ge3$。故一对根中恰有一个代表的平方已模 $2^{k+1}$ 同余于 $a$。对这个代表 $r$，$r$ 与 $r+2^k$ 给出两个提升，因为它们平方之差被 $2^{k+1}$ 整除；对另一个代表则两个提升都失败。两对根总共给出四根。任何高模数根都约化为低模数根，因此没有遗漏。
\end{proof}

\begin{example}[把奇素数幂与二的幂合并]
求 $x^2\equiv1\pmod{72}$ 的全部根类。
\end{example}
\begin{solve}
$72=8\cdot9$。模 $8$ 的根为 $1,3,5,7$，模 $9$ 的根为 $1,8$，所以共有 $4\cdot2=8$ 根。对模 $9$ 的第一类，列 $1+9t$（$0\le t<8$），其中奇数恰为 $1,19,37,55$；对第二类，列 $8+9t$，其中奇数恰为 $17,35,53,71$。它们模 $8$ 都是平方根 $1$，故全部根为 $1,17,19,35,37,53,55,71$。
\end{solve}

当 $a$ 与模数不互素时，根数还会变化。例如 $x^2\equiv0\pmod{p^k}$ 等价于 $p^{\lceil k/2\rceil}\mid x$，所以有 $p^{\lfloor k/2\rfloor}$ 个根类。若 $v_p(a)=t<k$ 且存在根，则 $t$ 必为偶数：两个模 $p^k$ 同余、其中一个素因子指数小于 $k$ 的整数具有相同的 $p$ 指数，故 $2v_p(x)=t$。写 $t=2s$、$x=p^sy$ 后，可以先约去 $p^{2s}$，归结为模 $p^{k-2s}$ 的可逆平方同余；最后恢复模 $p^k$ 的根时，还须计入不同的提升，不能只保留约分后的代表元。

\begin{corollary}
\label{cor:prime-three-mod-four}
若素数 $p\equiv3\pmod4$ 且 $p\mid x^2+y^2$，则 $p\mid x$ 且 $p\mid y$。
\end{corollary}
\begin{proof}
若 $p\nmid y$，则 $(xy^{-1})^2\equiv-1\pmod p$，与 $\left(\frac{-1}p\right)=-1$ 矛盾。故 $p\mid y$，进而 $p\mid x^2$，所以 $p\mid x$。
\end{proof}

\section{习题}

勒让德符号在奇素数模数下准确判定非零平方，雅可比符号在合数模数下只保留较粗的信息。处理一般模数时，先逐个检查素数幂，再用中国剩余定理合并；模二的幂与不可逆右端都需要单独分析。

\begin{enumerate}
 \item 列出模 $13$ 的非零平方剩余，并用欧拉判别法核对。
 \item 用互反律计算 $\left(\frac7{19}\right)$ 与 $\left(\frac{13}{17}\right)$。
 \item 求 $x^2\equiv4\pmod{45}$ 的全部解。
 \item 若奇素数 $p\mid a^2+1$，证明 $p\equiv1\pmod4$；解释为什么必须排除 $p=2$。
 \item\label{ex:quadratic-sign} 计算 $\left(\frac3{17}\right)$、$\left(\frac5{29}\right)$，注明每一步是否使用互反律的变号。
 \item\label{ex:quadratic-jacobi} 计算 $\left(\frac5{21}\right)$，判断 $x^2\equiv5\pmod{21}$ 是否可解，解释符号为 $1$ 为何仍可能无解。
 \item\label{ex:quadratic-two-power} 求 $x^2\equiv9\pmod{32}$ 的全部根类，并核对根数。
 \item\label{ex:quadratic-valuation} 求 $x^2\equiv9\pmod{27}$ 的全部根类，解释为什么约去 $9$ 后还要恢复多个模 $27$ 的类。
\end{enumerate}

\chapter{不定方程与平方和}
\label{ch:diophantine}

一个实数方程可能有一条连续曲线上的解，限制为整数后却可能没有解、只有有限多解，或形成特殊的无限序列。本章综合使用因式分解、同余、互素性与参数化。对于勾股数和平方和，我们不仅要给出漂亮的公式，还要说明为什么每个解都能被这些公式或判据覆盖。

\section{整数方程的解题框架}

整数方程的第一步是明确未知量范围。整除与因式分解常能把无限搜索化为有限个约数组合；同余可以给出无解障碍；不等式用于界定取值；参数化用于描述全部解。整数解必在每个模数下给出同余解，因此某个模数下无解足以否定整数解，但通过有限个模数的检验并不足以证明整数解存在。

一种实用的思考顺序是先检查奇偶和小模数，排除明显障碍；再寻找可分解的乘积或能提取的最大公因数；若仍有无限多个候选，就考虑参数化、递降或大小估计。这个顺序只是解题策略，不是保证所有方程都可解的算法。证明无解时找到一个矛盾即可，证明有解要构造对象，求全部解则必须同时控制所有可能。

\begin{example}
求 $xy=x+y+3$ 的全部正整数解。
\end{example}
\begin{solve}
改写为 $(x-1)(y-1)=4$。正整数解必须有 $x,y\ge2$，因此因子对只有 $(1,4),(2,2),(4,1)$，得到 $(x,y)=(2,5),(3,3),(5,2)$。因子枚举已覆盖所有可能，所以答案完整。
\end{solve}

\section{勾股数的参数化}

满足 $x^2+y^2=z^2$ 的正整数三元组称为\keyterm{勾股三元组（Pythagorean triple）}；若 $\gcd(x,y,z)=1$，称其本原（primitive）。本原解的三个数两两互素，否则某个公共素因子也会整除第三个数。

\begin{theorem}
\label{thm:pythagorean}
交换两条直角边后，每个本原勾股三元组唯一写为
\begin{equation}
 x=u^2-v^2,\qquad y=2uv,\qquad z=u^2+v^2,
 \label{eq:pythagorean-parameters}
\end{equation}
其中 $u>v>0$，$\gcd(u,v)=1$，且 $u,v$ 奇偶性不同；规定 $x$ 为奇边、$y$ 为偶边。
\end{theorem}
\begin{proof}
两边不能均偶，也不能均奇，否则平方和模 $4$ 余 $2$。故可令 $x,z$ 奇、$y$ 偶。由
\[
 \frac{z+x}{2}\frac{z-x}{2}=\left(\frac y2\right)^2
\]
知左侧两因子乘积为平方。它们互素：公因数整除其和 $z$ 与差 $x$。因此分别为 $u^2,v^2$，且 $u>v>0$，得到所示公式。互素性给出 $\gcd(u,v)=1$，$z$ 为奇数又给出二者奇偶性不同；$u^2=(z+x)/2,v^2=(z-x)/2$ 保证唯一性。

反之，公式直接满足勾股等式。若某个奇素数同时整除 $u^2-v^2$ 与 $u^2+v^2$，则它整除 $u,v$，矛盾；两者又都是奇数。因此所构造的三元组本原。
\end{proof}

一般正整数解由本原解同时乘以正整数得到。取 $(u,v)=(2,1),(3,2)$，分别得到 $(3,4,5)$ 与 $(5,12,13)$。

\begin{example}[参数条件怎样保证不漏解]
求所有斜边不超过 $20$ 的本原勾股三元组，不区分两直角边的次序。
\end{example}
\begin{solve}
由 $z=u^2+v^2\le20$ 且 $u>v\ge1$，有 $2\le u\le4$。$u=2$ 时只能取 $v=1$，得 $(3,4,5)$；$u=3$ 时 $v=1$ 与 $u$ 同奇，应排除，$v=2$ 得 $(5,12,13)$；$u=4$ 时 $v=1$ 得 $(15,8,17)$，$v=2$ 不互素，$v=3$ 又使斜边为 $25>20$。故只有这三组。若不要求本原，还要加入这些三元组的正整数倍，并重新检查斜边上界。
\end{solve}

参数 $u,v$ 的互素与奇偶条件不是装饰：例如 $(u,v)=(3,1)$ 给出 $(8,6,10)$，虽满足勾股等式，却不是本原解。证明参数化的反向时应检查“确实满足方程”和“确实满足本原要求”两件事，不能只代入平方等式就结束。

\section{无穷递降法}

\keyterm{无穷递降法（infinite descent）}的要点是：若一组解存在，就能构造同类且某个正整数参数严格更小的解；这与良序原理矛盾。

\begin{theorem}
\label{thm:fermat-four}
方程 $x^4+y^4=z^2$ 没有正整数解。
\end{theorem}
\begin{proof}
递降的目标是把任意正解变成右侧平方根更小的正解。之所以研究 $x^4+y^4=z^2$ 而不是直接研究 $x^4+y^4=z^4$，是因为两次勾股参数化会自然产生右侧为平方的同类方程，从而能够闭合递降过程。

若有解，取 $z$ 最小的一组。它满足 $\gcd(x,y)=1$：若 $d>1$ 同时整除 $x,y$，则 $d^4\mid z^2$，由唯一分解有 $d^2\mid z$，约去 $d^4$ 就得到右侧参数为 $z/d^2$ 的更小解。由奇偶性可令 $x$ 奇、$y$ 偶。本原勾股三元组 $(x^2,y^2,z)$ 给出
\[
 x^2=u^2-v^2,\qquad y^2=2uv,\qquad z=u^2+v^2,
\]
其中 $u,v$ 互素且奇偶不同。若 $u$ 偶而 $v$ 奇，则 $x^2\equiv3\pmod4$，不可能；所以 $u$ 奇、$v$ 偶。由 $(y/2)^2=u(v/2)$ 及两因子互素，得 $u=s^2,v=2t^2$，于是
\[
 x^2+(2t^2)^2=s^4.
\]
这是本原勾股三元组：若素数同时整除 $x,t$，由 $x^2=s^4-4t^4$ 就也整除 $s$，与 $\gcd(u,v)=1$ 矛盾。再次参数化得
\[
 x=r^2-w^2,\qquad t^2=rw,\qquad s^2=r^2+w^2,
\]
其中 $r,w$ 互素且为正整数。乘积为平方迫使 $r=a^2,w=b^2$，因此 $a^4+b^4=s^2$。而 $s^2=u<z$，特别有 $s<z$，构造出了更小的正整数解，矛盾。
\end{proof}

因此 $x^4+y^4=z^4$ 也无正整数解，只需把右侧视为 $(z^2)^2$。这一论证只处理指数 $4$，并未证明所有指数的情形。

回看递降证明，需要逐项核对新构造的数仍为整数、仍为正数、仍满足同一种方程，并且选定参数严格变小。仅把数值“变小一些”或把方程变成另一类方程，都不足以产生矛盾。这里新解是 $(a,b,s)$，缩小的是右侧平方根，从 $z$ 变为 $s$；它恰好符合最初选择最小 $z$ 的标准。

\section{两平方和定理}

\begin{lemma}
\label{lem:prime-two-squares}
奇素数 $p$ 能写成两个整数的平方和，当且仅当 $p\equiv1\pmod4$。
\end{lemma}
\begin{proof}
必要性来自平方模 $4$ 的余数：奇数平方和只能由一奇一偶产生，因而余 $1$。

若 $p\equiv1\pmod4$，由式\eqref{eq:quadratic-supplements}可取 $t$ 使 $t^2\equiv-1\pmod p$。令 $h=\lfloor\sqrt p\rfloor$，考虑 $(h+1)^2>p$ 个整数对 $(u,v)$，其中 $0\le u,v\le h$，将它们按 $u-tv\pmod p$ 分入 $p$ 个盒子。两个不同的对有相同余数，作差得到不全为零的 $x,y$，满足 $|x|,|y|\le h$ 且 $x\equiv ty\pmod p$。于是
\[
 0<x^2+y^2\le2h^2<2p,\qquad p\mid x^2+y^2.
\]
只能有 $x^2+y^2=p$。
\end{proof}

\begin{theorem}[两平方和定理（sum of two squares theorem）]
\label{thm:two-squares}
正整数 $n$ 可写成 $x^2+y^2$（$x,y\in\mathbb Z$），当且仅当每个满足 $p\equiv3\pmod4$ 的素因子 $p$ 在 $n$ 的分解中具有偶数指数。允许 $x$ 或 $y$ 为零。
\end{theorem}
\begin{proof}
若这样的 $p$ 整除 $x^2+y^2$，由推论\ref{cor:prime-three-mod-four}，$p$ 同时整除 $x,y$，所以可从和式提出 $p^2$。反复进行，$p$ 的指数只能为偶数。

反之，$2=1^2+1^2$，模 $4$ 余 $1$ 的素数由引理\ref{lem:prime-two-squares}表示为两平方和，模 $4$ 余 $3$ 的素数的偶次幂本身为平方。最后使用恒等式
\begin{equation}
 (a^2+b^2)(c^2+d^2)=(ac-bd)^2+(ad+bc)^2
 \label{eq:two-square-product}
\end{equation}
逐个相乘即可得到 $n$ 的表示。
\end{proof}

例如 $65=5\cdot13$ 满足判据，并有 $65=1^2+8^2=4^2+7^2$；$21=3\cdot7$ 不满足判据，因而无表示。判据解决可表示性，完整枚举全部表示则是进一步的问题。

\begin{example}[从判据走向构造]
判断 $325$ 是否为两平方和，并由较小数的表示构造两组表示。
\end{example}
\begin{solve}
$325=5^2\cdot13$，没有模 $4$ 余 $3$ 的素因子，故可以表示。先用 $25=3^2+4^2$、$13=2^2+3^2$，代入式\eqref{eq:two-square-product}得
\[
 325=(3\cdot2-4\cdot3)^2+(3\cdot3+4\cdot2)^2=6^2+17^2.
\]
把第二个表示改为 $13=2^2+(-3)^2$，得到 $325=18^2+1^2$。还可用 $25=5^2+0^2$ 得到 $325=10^2+15^2$。这些构造验证了多组表示的存在；若要求全部表示，还应加上枚举范围或其他完备性论证。
\end{solve}

\section{四平方和专题（选学）}

\begin{theorem}[拉格朗日四平方和定理；此处不证]
每个非负整数都能写成四个整数的平方和。
\end{theorem}

该结论称为\keyterm{四平方和定理（four-square theorem）}，其完整证明可见参考资料\cite{mit-four-squares}。与两平方和相比，四个变量消除了可表示性的限制。例如 $7=2^2+1^2+1^2+1^2$，但它不是两平方和；它也不是三平方和，因为平方模 $8$ 只可能余 $0,1,4$，三个这样的余数不可能相加得到 $7$。这里引用的是存在性结论，并未给出高效求表示的算法。

\section{习题}

求整数方程的全部解，需要把每次变换的正反方向都说明。勾股参数化依赖互素乘积为平方时各因子为平方；递降依赖正整数不能无限下降；两平方和判据则将整体可表示性化为素因子的类型与指数条件。

\begin{enumerate}
 \item 求斜边为 $25$ 的全部正整数勾股三元组，不区分两条直角边的次序。
 \item 判断 $45,75,130$ 哪些可表示为两个整数的平方和，并给出可表示者的一组表示。
 \item 证明不存在面积为奇数的整数边长直角三角形。
 \item 求 $xy=2x+3y$ 的全部正整数解，并说明枚举为何完整。
 \item\label{ex:dio-factors} 求 $x^2-y^2=45$ 的全部非负整数解。
 \item\label{ex:dio-hypotenuse} 证明本原勾股三元组的斜边的每个素因子都模 $4$ 余 $1$。
 \item\label{ex:dio-two-square-list} 判断 $50,77,98,325$ 是否为两平方和，并给出可表示者的一组表示。
 \item\label{ex:dio-area} 证明整数边长直角三角形的面积一定被 $6$ 整除。提示：先处理本原三元组，再乘回公因数。
\end{enumerate}

\chapter{连分数与有理逼近}
\label{ch:continued-fractions}

小数展开按十进制位逼近实数，连分数则按分母的大小寻找有理近似。例如 $\sqrt2$ 的近似 $1,3/2,7/5,17/12$，分母并不是 $10$ 的幂，却能迅速提高精度。连分数算法其实就是不断“取整数部分，再把剩余部分取倒数”，与欧几里得算法紧密相连。本章先处理有限展开，再建立无穷展开的误差估计，最后为 Pell 方程准备平方根的周期结构。

\section{有限简单连分数与有理数}

\begin{definition}
\keyterm{简单连分数（simple continued fraction）}写为
\begin{equation}
 [a_0;a_1,\ldots,a_n]
 =a_0+\cfrac1{a_1+\cfrac1{\ddots+\cfrac1{a_n}}},
 \label{eq:finite-continued-fraction}
\end{equation}
其中 $a_0\in\mathbb Z$，$a_1,\ldots,a_n$ 为正整数；$n=0$ 时值为 $a_0$。
\end{definition}

对有理数 $\alpha$，令 $a_0=\lfloor\alpha\rfloor$。若余下的小数部分非零，取倒数再重复取整；这是欧几里得算法的另一种写法，分母随带余除法下降，因而有限步终止。反之，有限连分数由有理运算构成，值必为有理数。

若末项为 $1$，可用 $[\ldots,a,1]=[\ldots,a+1]$ 合并。规定非整数的有限展开末项至少为 $2$，则每步的整数部分被唯一确定，得到唯一的规范展开。例如
\[
 415=4\cdot93+43,\quad93=2\cdot43+7,\quad43=6\cdot7+1,
\]
故 $415/93=[4;2,6,7]$。

将除法写成分数形式，更容易看见每个部分商的来源：
\[
 \frac{415}{93}=4+\frac{43}{93},\qquad
 \frac{93}{43}=2+\frac7{43},\qquad
 \frac{43}{7}=6+\frac17.
\]
依次代回就得到 $4+1/(2+1/(6+1/7))$。规范条件也有实际作用，例如 $[1;2]=3/2=[1;1,1]$，所以不限制最后一项就谈不上展开唯一。对负数仍采用向下取整：$-7/3=-3+2/3$，从而规范展开为 $[-3;1,2]$。

\section{渐近分数与递推关系}

截取前 $k+1$ 项得到的 $P_k/Q_k=[a_0;a_1,\ldots,a_k]$ 称为\keyterm{渐近分数（convergent，又称收敛分数）}。令
\begin{equation}
 P_{-2}=0,\quad P_{-1}=1,\qquad Q_{-2}=1,\quad Q_{-1}=0,
 \label{eq:convergent-initial}
\end{equation}
并递推
\begin{equation}
 P_k=a_kP_{k-1}+P_{k-2},\qquad
 Q_k=a_kQ_{k-1}+Q_{k-2}.
 \label{eq:convergent-recurrence}
\end{equation}
这些公式可由矩阵恒等式直接验证：
\[
 \begin{pmatrix}a_0&1\\1&0\end{pmatrix}\cdots
 \begin{pmatrix}a_k&1\\1&0\end{pmatrix}
 =\begin{pmatrix}P_k&P_{k-1}\\Q_k&Q_{k-1}\end{pmatrix}.
\]
每个左侧矩阵表示分式变换 $t\mapsto a_j+1/t$，连乘表示依次代入；取行列式得到
\begin{equation}
 P_kQ_{k-1}-P_{k-1}Q_k=(-1)^{k+1},\qquad k\ge0.
 \label{eq:convergent-determinant}
\end{equation}
因此 $\gcd(P_k,Q_k)=1$，相邻渐近分数之差的绝对值为 $1/(Q_kQ_{k-1})$（$k\ge1$）。$Q_0=1$，$Q_k>0$，且 $Q_{k+2}\ge Q_{k+1}+Q_k$，所以无限展开的分母趋于无穷。

初始值中的负下标只是统一递推的辅助记号，不代表额外的连分数项。代入 $k=0$ 得 $P_0=a_0,Q_0=1$；代入 $k=1$ 得 $P_1=a_1a_0+1,Q_1=a_1$，恰为 $a_0+1/a_1$。对 $[4;2,6,7]$，计算如下：
\[
\begin{array}{c|r|r|r|c}
k&a_k&P_k&Q_k&P_k/Q_k\\\hline
0&4&4&1&4\\
1&2&9&2&9/2\\
2&6&58&13&58/13\\
3&7&415&93&415/93
\end{array}
\]
例如 $58\cdot2-9\cdot13=-1$，所以 $58,13$ 的任何公因数都整除 $1$。行列式恒等式同时解释了约分后的分母为什么已经是 $Q_k$，以及相邻分数为何彼此接近。

\section{无限连分数与误差估计}

对无理实数 $\alpha$，定义完整商 $\alpha_0=\alpha$、$a_k=\lfloor\alpha_k\rfloor$、$\alpha_{k+1}=1/(\alpha_k-a_k)$。过程不会终止，否则反向代入会使 $\alpha$ 为有理数；并且 $k\ge1$ 时 $\alpha_k>1$、$a_k\ge1$。矩阵恒等式给出
\begin{equation}
 \alpha=\frac{P_k\alpha_{k+1}+P_{k-1}}
 {Q_k\alpha_{k+1}+Q_{k-1}}.
 \label{eq:complete-quotient}
\end{equation}
由式\eqref{eq:convergent-determinant}，
\begin{equation}
 \alpha-\frac{P_k}{Q_k}
 =\frac{(-1)^k}{Q_k(Q_k\alpha_{k+1}+Q_{k-1})},\qquad
 \left|\alpha-\frac{P_k}{Q_k}\right|<\frac1{Q_kQ_{k+1}}.
 \label{eq:convergent-error}
\end{equation}
后一不等式使用 $\alpha_{k+1}>a_{k+1}$。因此偶数下标的渐近分数在 $\alpha$ 左侧，奇数下标的在右侧，且都收敛于 $\alpha$。

误差式也说明下一个部分商的意义：若 $a_{k+1}$ 较大，则 $Q_{k+1}=a_{k+1}Q_k+Q_{k-1}$ 较大，当前近似 $P_k/Q_k$ 的误差就更小。误差应与分母一起比较，不能只看小数位数；分母很大的分数获得较小误差，本身未必说明它是高效的逼近。

\begin{example}[不依赖小数展开的误差计算]
比较 $7/5$、$17/12$ 与 $\sqrt2$ 的大小，并估计后一个近似的误差。
\end{example}
\begin{solve}
因为 $7^2-2\cdot5^2=-1$、$17^2-2\cdot12^2=1$，有 $7/5<\sqrt2<17/12$。有理化给出精确表达式
\[
 \frac{17}{12}-\sqrt2
 =\frac{1}{12^2(17/12+\sqrt2)}.
\]
下一渐近分数为 $41/29$，所以式\eqref{eq:convergent-error}还给出误差小于 $1/(12\cdot29)$。前式来自平方差，后式来自连分数的一般理论；两条路线相互核对，并为后章的 Pell 方程建立联系。
\end{solve}

反过来，任意满足部分商条件的无限序列也确定一个无理极限。由递推，偶数子列严格递增、奇数子列严格递减，二者交错且相邻差趋于零，所以趋于同一实数。若极限为既约分数 $A/B$，它不同于每个渐近分数，故误差至少为 $1/(BQ_k)$；但嵌套区间给出的误差至多为 $1/(Q_kQ_{k+1})$，当 $Q_{k+1}>B$ 时矛盾。分式变换的连续性又使极限的取整展开恢复原部分商序列。

\begin{lemma}[第二类逼近性质]
\label{lem:best-approximation}
设 $\alpha$ 无理，$k\ge0$。若整数 $a,b$ 满足 $0<b<Q_{k+1}$，且 $a/b\ne P_k/Q_k$，则
\[
 |b\alpha-a|\ge|Q_k\alpha-P_k|.
\]
\end{lemma}
\begin{proof}
由式\eqref{eq:convergent-determinant}，相邻两列组成的矩阵行列式为 $\pm1$，消元所得逆矩阵仍为整数矩阵，故唯一存在整数 $u,v$ 使
\[
 (a,b)=u(P_k,Q_k)+v(P_{k+1},Q_{k+1}).
\]
$v=0$ 会使分数相等，故 $v\ne0$。由 $0<b<Q_{k+1}$，若 $v>0$ 必有 $u<0$，若 $v<0$ 必有 $u>0$。两误差 $Q_k\alpha-P_k$ 与 $Q_{k+1}\alpha-P_{k+1}$ 符号相反，因此乘以异号的 $u,v$ 后同号相加，绝对值至少为 $|Q_k\alpha-P_k|$。
\end{proof}

\begin{theorem}[勒让德逼近判据（Legendre's approximation criterion）]
\label{thm:legendre-approximation}
若 $a/b$ 既约、$b>0$，且 $|\alpha-a/b|<1/(2b^2)$，则 $a/b$ 是无理数 $\alpha$ 的一个渐近分数。
\end{theorem}
\begin{proof}
选取最大的 $k$ 使 $Q_k\le b$，则 $b<Q_{k+1}$。若分数不等于 $P_k/Q_k$，由引理\ref{lem:best-approximation}得
\[
 1\le|aQ_k-bP_k|
 \le Q_k|a-b\alpha|+b|Q_k\alpha-P_k|
 \le(Q_k+b)|a-b\alpha|<1,
\]
矛盾。这里最后使用 $Q_k\le b$ 与假定误差界。
\end{proof}

该引理比较的是 $|b\alpha-a|$，不是在任意分母范围内无条件最小化 $|\alpha-a/b|$；两类“最佳逼近”不能混用。

勒让德判据的方向同样需要留意：满足 $1/(2b^2)$ 的严格误差界，足以保证一个既约分数是渐近分数；它没有断言每个渐近分数都满足这个更强的误差界。下一章会从整数方程导出这一误差界，进而把无限多个可能的分数缩小为一个按递推生成的序列。

\section{周期连分数与二次无理数}

\keyterm{二次无理数（quadratic irrational）}是满足整系数二次方程的无理实数。以下证明后章所需的平方根周期性。

\begin{theorem}[平方根的周期展开]
\label{thm:sqrt-period}
若 $D$ 为非平方正整数，$a_0=\lfloor\sqrt D\rfloor$，则
\[
 \sqrt D=[a_0;\overline{a_1,\ldots,a_{\ell-1},2a_0}],
\]
其中横线表示不断重复，$\ell\ge1$ 为最小周期长度。
\end{theorem}
\begin{proof}
将完整商写为 $\alpha_j=(\sqrt D+B_j)/C_j$，从 $B_0=0,C_0=1$ 开始，直接有理化得到
\begin{equation}
 B_{j+1}=a_jC_j-B_j,\qquad
 C_{j+1}=\frac{D-B_{j+1}^2}{C_j}.
 \label{eq:surd-recurrence}
\end{equation}
由 $C_j\mid D-B_j^2$ 归纳知各 $B_j,C_j$ 均为整数且保持这一整除条件。

记 $\alpha_j'=(B_j-\sqrt D)/C_j$ 为共轭值。$\alpha_1>1$ 且 $-1<\alpha_1'=1/(-\sqrt D-a_0)<0$。若此性质对 $j$ 成立，则 $\alpha_{j+1}>1$，且 $\alpha_{j+1}'=1/(\alpha_j'-a_j)\in(-1,0)$。所以对 $j\ge1$，$C_j>0$，并由
\[
 \alpha_j-\alpha_j'=\frac{2\sqrt D}{C_j}>1
\]
得 $C_j<2\sqrt D$，又有 $-\sqrt D<B_j<\sqrt D$。整数状态 $(B_j,C_j)$ 只有有限种，故必重复。

转移在这些状态上可逆：给定 $\alpha_{j+1}'$，由 $-1/\alpha_{j+1}'=a_j-\alpha_j'\in(a_j,a_j+1)$ 唯一恢复 $a_j$，进而恢复 $\alpha_j$。因此从 $\alpha_1$ 起即为一个循环。数 $\beta=\sqrt D+a_0$ 同样满足 $\beta>1,-1<\beta'<0$，且其下一完整商为 $\alpha_1$，所以它是循环中 $\alpha_1$ 的前一项，即 $\alpha_\ell=\beta$。其整数部分为 $2a_0$，定理得证。
\end{proof}

证明还给出整数算法：用式\eqref{eq:surd-recurrence}迭代，并取 $a_j=\lfloor(a_0+B_j)/C_j\rfloor$。无需浮点近似 $\sqrt D$，因为 $0<\sqrt D-a_0<1$ 不会改变该商的整数部分。对 $j\ge1$，$C_j=1$ 当且仅当 $B_j=a_0$，也即当且仅当 $\ell\mid j$；这由 $-1<B_j-\sqrt D<0$ 和循环中状态互异得到。

证明里“状态有限”首先只给出最终重复，而“转移可逆”才进一步排除进入循环前的一段尾巴。若一个确定过程允许两个不同状态进入同一状态，就可能先经过非周期部分再进入周期。这里能从后一个完整商的共轭唯一恢复前一步，所以从 $\alpha_1$ 起就是循环。这是周期证明中不能省略的逻辑环节。

\begin{example}[平方根展开的整数计算表]
求 $\sqrt{14}$ 的周期连分数。
\end{example}
\begin{solve}
$a_0=3$，从 $(B_0,C_0)=(0,1)$ 开始，使用式\eqref{eq:surd-recurrence}，得到
\[
\begin{array}{c|r|r|r}
j&B_j&C_j&a_j\\\hline
0&0&1&3\\
1&3&5&1\\
2&2&2&2\\
3&2&5&1\\
4&3&1&6\\
5&3&5&1
\end{array}
\]
第 $5$ 行的状态首次回到第 $1$ 行，故 $\sqrt{14}=[3;\overline{1,2,1,6}]$，最小周期为 $4$。例如第二次更新是 $B_2=1\cdot5-3=2$、$C_2=(14-2^2)/5=2$。整个过程只使用整数除法，避免了舍入误差导致部分商取错的问题。
\end{solve}

一般的拉格朗日周期定理（Lagrange's periodicity theorem）断言：实数的无限简单连分数最终周期，当且仅当它为二次无理数。这里仅引用一般情形，证明见\cite{mit-periodic}；“最终周期”允许周期前有有限项，不能与从起点开始的纯周期混同。后章只依赖上面已证明的平方根情形。

\section{习题}

连分数递推同时控制分子、分母和误差，相邻渐近分数的行列式为正负一是核心恒等式。平方根展开的周期来自有限整数状态及可逆转移；计算中应分清部分商、完整商和渐近分数三个不同对象。

\begin{enumerate}
 \item 由递推计算 $[4;2,6,7]$ 的所有渐近分数。
 \item 证明 $\sqrt2=[1;\overline2]$，计算前五个渐近分数并核对误差符号。
 \item 用整数状态递推求 $\sqrt3$ 与 $\sqrt{13}$ 的周期。
 \item 若 $a_k=1$ 对所有 $k\ge1$ 成立，说明分母递推与斐波那契数列的关系。
 \item\label{ex:cf-negative} 求 $-7/3$ 的规范简单连分数，说明第一项为什么是 $-3$ 而不是 $-2$。
 \item\label{ex:cf-sqrt-seven} 求 $\sqrt7$ 的最小周期与前四个渐近分数，并核对它们在 $\sqrt7$ 两侧交错。
 \item\label{ex:cf-neighbors} 设正整数 $a,b,c,d$ 满足 $a/b<c/d$ 且 $bc-ad=1$。若整数 $u,v$ 满足 $v>0$、$a/b<u/v<c/d$，证明 $v\ge b+d$。
 \item\label{ex:cf-golden} 设 $\alpha=[1;\overline1]$，求 $\alpha$ 满足的二次方程，并由 $\alpha>1$ 确定其值。
\end{enumerate}

\chapter{Pell 方程}
\label{ch:pell}

方程 $x^2-2y^2=1$ 的解 $(3,2),(17,12),(99,70)$ 很快变大，逐个试代 $y$ 不便于解释其规律。前章已经说明，连分数能够产生非常接近平方根的有理数；本章利用 $x/y$ 对 $\sqrt D$ 的近似，把解的搜索化为渐近分数的检查，再用带根号数的乘法描述全部解。

\section{方程形式与平方模数的区别}

固定非平方正整数 $D$，\keyterm{佩尔方程（Pell equation）}为
\begin{equation}
 x^2-Dy^2=1,\qquad x,y\in\mathbb Z.
 \label{eq:pell-equation}
\end{equation}
$y=0$ 给出平凡解 $(\pm1,0)$；非平凡解可先取 $x,y>0$，再恢复符号。若 $D=s^2$ 是平方，则 $(x-sy)(x+sy)=1$，整数因子只能同为 $1$ 或同为 $-1$，因而只有平凡解。非平方条件实质上改变了方程的结构。

对形如 $u=x+y\sqrt D$ 的数定义共轭 $u'=x-y\sqrt D$ 与\keyterm{范数（norm）} $N(u)=uu'=x^2-Dy^2$。由直接乘法可知 $N(uv)=N(u)N(v)$。整数系数在乘法下保持；若 $N(u)=1$，其逆元就是共轭，仍有整数系数。

这里的“范数”不是长度，也不要求为非负数；它是一个便于记录方程左侧的代数表达式。若 $u=x+y\sqrt D$、$v=z+w\sqrt D$，则
\[
 uv=(xz+Dyw)+(xw+yz)\sqrt D,
\]
所以整数系数在乘法下保持的事实完全来自普通展开。例如 $3+2\sqrt2$ 的范数是 $1$，平方得到 $17+12\sqrt2$，其范数仍为 $1$，故自动产生另一组解。

\section{连分数方法与基本解}

设 $P_k/Q_k$ 为 $\sqrt D$ 的渐近分数，$B_j,C_j$ 为式\eqref{eq:surd-recurrence}的状态。

\begin{lemma}
\label{lem:pell-convergent-norm}
对每个 $k\ge0$，
\begin{equation}
 P_k^2-DQ_k^2=(-1)^{k+1}C_{k+1}.
 \label{eq:pell-convergent-norm}
\end{equation}
而任意满足 $x,y>0$、$x^2-Dy^2=\pm1$ 的解，$x/y$ 必为 $\sqrt D$ 的渐近分数。
\end{lemma}
\begin{proof}
将 $\alpha_{k+1}=(\sqrt D+B_{k+1})/C_{k+1}$ 代入式\eqref{eq:complete-quotient}，比较有理部分和 $\sqrt D$ 的系数，得
\[
 P_k=B_{k+1}Q_k+C_{k+1}Q_{k-1},\qquad
 DQ_k=B_{k+1}P_k+C_{k+1}P_{k-1}.
\]
消去 $B_{k+1}$，再用行列式恒等式即得范数公式。

若 $x^2-Dy^2=\pm1$，则 $\gcd(x,y)=1$，并有
\[
 \left|\frac xy-\sqrt D\right|
 =\frac1{y^2(x/y+\sqrt D)}<\frac1{2y^2}.
\]
这里 $D\ge2$，且 $(x/y)^2=D\pm1/y^2\ge D-1\ge1$，所以 $x/y\ge1$，分母括号大于 $2$。由定理\ref{thm:legendre-approximation}，$x/y$ 必是渐近分数。
\end{proof}

\begin{theorem}[基本解的周期公式]
\label{thm:pell-fundamental}
设 $\sqrt D$ 的最小周期长度为 $\ell$。方程\eqref{eq:pell-equation}存在正整数解，其中 $x$ 最小的一组称为\keyterm{基本解（fundamental solution）}，并由
\[
 (x_1,y_1)=
 \begin{cases}
 (P_{\ell-1},Q_{\ell-1}),&\ell\text{ 为偶数},\\
 (P_{2\ell-1},Q_{2\ell-1}),&\ell\text{ 为奇数}
 \end{cases}
\]
给出。
\end{theorem}
\begin{proof}
由定理\ref{thm:sqrt-period}之后的状态结论，$C_j=1$ 当且仅当 $\ell\mid j$。因此渐近分数给出范数 $\pm1$ 当且仅当 $k+1=r\ell$，相应范数为 $(-1)^{r\ell}$。最早的正范数出现于所列下标，这也证明了解的存在。任意正解都由引理\ref{lem:pell-convergent-norm}来自渐近分数，而这些分子的值严格递增，所以最早出现者即为基本解。
\end{proof}

计算时先求周期，再用整数递推计算相应的 $P_k,Q_k$，最后代回原方程核查。也可以依次测试渐近分数的范数直到遇到 $1$；上述定理保证过程终止。

\begin{example}[从一个周期中读出基本解]
求 $x^2-14y^2=1$ 的基本解。
\end{example}
\begin{solve}
前章已算出 $\sqrt{14}=[3;\overline{1,2,1,6}]$，周期 $\ell=4$ 为偶数，所以取下标 $3$ 的渐近分数。前四个渐近分数依次为 $3,4,11/3,15/4$；对应的 $P_k^2-14Q_k^2$ 依次为 $-5,2,-5,1$。因此基本解为 $(15,4)$，代回得 $225-14\cdot16=1$。此处下标从 $0$ 开始，所以“下标 $\ell-1$”意味着取前 $\ell$ 个部分商，而不是把周期末项 $6$ 也算进去。
\end{solve}

若周期为奇数，第一周期末得到范数 $-1$，必须到第二周期对应位置才得到范数 $1$。例如 $\sqrt5=[2;\overline4]$，第一项 $2/1$ 给出 $2^2-5=-1$，再取一项得到 $9/4$，才有 $9^2-5\cdot4^2=1$。这个符号差别解释了基本解公式为何按周期奇偶分成两种情况。

\section{全部正整数解的生成}

\begin{theorem}
\label{thm:pell-all}
若 $(x_1,y_1)$ 是基本解，则全部正整数解恰由
\begin{equation}
 x_n+y_n\sqrt D=(x_1+y_1\sqrt D)^n,
 \qquad n=1,2,\ldots
 \label{eq:pell-all-solutions}
\end{equation}
给出。
\end{theorem}
\begin{proof}
令 $U=x_1+y_1\sqrt D>1$。正整数幂的系数为正整数，范数为 $1$，所以确实给出解。

反之，任取正解对应的 $V=x+y\sqrt D>1$。因为 $U>1$，可选整数 $n\ge0$ 使 $U^n\le V<U^{n+1}$。令 $W=VU^{-n}$，由 $U^{-1}=x_1-y_1\sqrt D$，$W$ 有整数系数、范数为 $1$，且 $1\le W<U$。写 $W=X+Y\sqrt D$，则
\[
 X=\frac{W+W^{-1}}2,\qquad
 Y=\frac{W-W^{-1}}{2\sqrt D}.
\]
若 $W>1$，则 $X,Y>0$，且 $X<x_1$：函数 $(t+t^{-1})/2$ 在 $t>1$ 上严格递增，可由两值作差验证。这与基本解的最小性矛盾。故 $W=1$，即 $V=U^n$；因 $V>1$，必有 $n\ge1$。
\end{proof}

展开乘法得到不含根号的递推
\[
 x_{n+1}=x_1x_n+Dy_1y_n,\qquad
 y_{n+1}=y_1x_n+x_1y_n.
\]
全部整数解还包括 $(\pm1,0)$ 以及对每个正解独立改变两个坐标的符号。

定理的前半部分只说明“取幂总能造出解”，后半部分才证明“每个解都是这种幂”。后半部分的思想是用基本单位 $U$ 反复缩小任意解 $V$，把它移入区间 $[1,U)$；若得到的数严格大于 $1$，它就对应比基本解还小的正解，产生矛盾。因此得到的只能是 $1$，原来的 $V$ 必为 $U$ 的整数幂。这比验证前几组解都符合递推强得多。

\begin{example}
求 $x^2-2y^2=1$ 的前三组正整数解。
\end{example}
\begin{solve}
$\sqrt2=[1;\overline2]$，周期长为 $1$，基本解取下标 $1$ 的渐近分数 $3/2$，故为 $(3,2)$。再由 $(3+2\sqrt2)^2=17+12\sqrt2$、$(3+2\sqrt2)^3=99+70\sqrt2$，得到 $(17,12)$、$(99,70)$。定理\ref{thm:pell-all}保证继续取幂可得到全部正解。
\end{solve}

\section{负 Pell 方程与相关问题}

\begin{proposition}
\label{prop:negative-pell}
$x^2-Dy^2=-1$ 有整数解，当且仅当 $\sqrt D$ 的最小周期长度 $\ell$ 为奇数。可解时最小正解为 $(P_{\ell-1},Q_{\ell-1})$。
\end{proposition}
\begin{proof}
任意正解均来自渐近分数；由式\eqref{eq:pell-convergent-norm}，范数为 $-1$ 等价于 $k+1=r\ell$ 且 $r\ell$ 为奇数。这在 $\ell$ 为奇数时且仅此时可发生，最早下标为 $\ell-1$。若存在任意整数解，取坐标绝对值即可得到正解，因为 $D\ge2$ 排除了零坐标。
\end{proof}

若最小负范数正解对应 $V_0=x_0+y_0\sqrt D$，正 Pell 方程基本解对应 $U$，则负方程全部正解为 $V_0U^n$（$n\ge0$）。确实，任一其他负范数正解 $V\ge V_0$ 的商 $V/V_0=-VV_0'$ 有整数系数、范数 $1$ 且不小于 $1$，所以由定理\ref{thm:pell-all}是 $U$ 的非负整数幂。

例如 $D=2$ 时 $1^2-2\cdot1^2=-1$；$D=3$ 的周期长为 $2$，负方程无解，也可用模 $3$ 的平方余数直接排除。一般方程 $x^2-Dy^2=N$ 的解乘以范数 $1$ 的元素仍是解，但可能存在多个互不相连的解族，不能直接照搬单一基本解生成全部解的结论。

\section{从整数问题识别 Pell 方程}

\begin{example}[相邻整数的平方关系]
求满足 $m^2=2n(n+1)+1$ 的非负整数 $n$ 的生成方法，其中 $m$ 为正整数。
\end{example}
\begin{solve}
令 $X=2n+1$，原式等价于 $X^2-2m^2=-1$。负 Pell 方程的全部正解满足
\[
 X_j+m_j\sqrt2=(1+\sqrt2)(3+2\sqrt2)^j,\qquad j\ge0.
\]
这些解中的 $X_j$ 都是奇数，因为 $X_j^2=2m_j^2-1$ 为奇数。因此 $n_j=(X_j-1)/2$ 是非负整数，反向代入也成立。前三组 $(X_j,m_j)$ 为 $(1,1),(7,5),(41,29)$，得到 $n=0,3,20$。两次变换都可逆，且奇偶条件得到验证，所以这一生成方法覆盖全部所求解。
\end{solve}

识别 Pell 型方程时通常先配方，使某个一次表达式的平方减去另一平方的固定倍数等于常数；随后检查常数是否为 $1$ 或 $-1$，并检查新变量是否还受奇偶性或整除条件限制。若常数不是 $\pm1$，本章的范数乘法仍能生成解，但完备性需要另作讨论。

\section{习题}

连分数保证基本解存在并给出寻找方法，范数乘法保证取幂生成新解，最小性论证保证不遗漏任何正解。应用中必须检查配方变换是否可逆，以及新变量是否还带有奇偶或整除约束。

\begin{enumerate}
 \item 求 $D=3,5,13$ 时正 Pell 方程的基本解，并代入验证。
 \item 写出 $x^2-2y^2=-1$ 的前三组正解。
 \item 证明三角数 $n(n+1)/2$ 为平方 $m^2$ 等价于 $(2n+1)^2-8m^2=1$，求前两个正整数 $n$。
 \item 由整数递推证明正 Pell 方程的正解有无穷多组，且两坐标均严格递增。
 \item\label{ex:pell-seven} 求 $x^2-7y^2=1$ 的基本解和第二组正解。
 \item\label{ex:pell-negative-five} 求 $x^2-5y^2=-1$ 的前三组正解，并核对范数。
 \item\label{ex:pell-second-order} 设 $(x_n,y_n)$ 由基本解 $(x_1,y_1)$ 生成，补充 $(x_0,y_0)=(1,0)$，证明 $x_{n+2}=2x_1x_{n+1}-x_n$、$y_{n+2}=2x_1y_{n+1}-y_n$。
 \item\label{ex:pell-obstruction} 证明若 $D$ 有一个模 $4$ 余 $3$ 的素因子，则负 Pell 方程无整数解；说明这是一条必要条件，不能仅凭排除此障碍就宣布有解。
\end{enumerate}

\chapter{素数分布初步（选学）}
\label{ch:prime-distribution}

本章先用整除与同余研究特殊数列，再讨论正项级数与渐近式；后两部分使用级数收敛和实函数极限的基本概念。

“素数有无穷多个”只说明没有最后一个素数，并不说明素数出现得有多频繁。要比较稀疏程度，可以数出某个区间的素数，也可以考察素数倒数的累积增长。本章把初等可证的结论和需要进一步理论的定理明确分开，使读者理解每种说法究竟提供多少信息。

\section{特殊等差数列中的素数}

\begin{proposition}
形如 $4k+3$ 的素数和形如 $4k+1$ 的素数都各有无穷多个。
\end{proposition}
\begin{proof}
若模 $4$ 余 $3$ 的素数只有 $p_1,\ldots,p_r$，令 $N=4p_1\cdots p_r-1$。$N$ 是模 $4$ 余 $3$ 的奇数，故其素因子中至少一个 $q$ 模 $4$ 余 $3$，否则素因子全部余 $1$，乘积也余 $1$。但每个 $p_i$ 都不能整除 $N$，矛盾。

若模 $4$ 余 $1$ 的素数只有 $q_1,\ldots,q_s$，令 $A=2q_1\cdots q_s$，取奇数 $A^2+1$ 的素因子 $q$。此时 $q\nmid A$ 且 $A^2\equiv-1\pmod q$，由补充公式知 $q\equiv1\pmod4$；它又不同于原有全部 $q_i$，矛盾。空列表按空乘积为 $1$ 处理，同样成立。
\end{proof}

更一般的\keyterm{狄利克雷素数定理（Dirichlet's theorem on primes in arithmetic progressions）}断言：若正整数 $a,m$ 互素，则 $a+km$（$k\ge0$）含无穷多个素数。互素条件不可删去，因为公共因子会整除数列中的每一项。这里仅陈述一般定理，其解析证明见\cite{kedlaya-dirichlet}。

例如 $6k+5$ 与 $6k+1$ 都包含无穷多个素数；但 $6k+3$ 中除 $3$ 本身外全是合数。若 $d=\gcd(a,m)>1$，数列中的素数项只能等于 $d$，且只有在 $d$ 本身为素数时才有这种可能，因此至多出现一个素数项。反过来，互素条件保证无穷多个素数，是一个远强于“没有固定公因子障碍”的定理，不能仅凭障碍消失就宣布结论成立。

\section{素数间隔与素数倒数和}

给定 $L\ge1$，令 $N=(L+1)!$，则 $N+2,N+3,\ldots,N+L+1$ 都是合数：第 $j$ 项被 $j$ 整除且严格大于 $j$。这构造了任意长的连续合数段，但没有断言空隙长度与其所处位置成正比。

\begin{theorem}[素数倒数和发散]
$\sum_{p\text{ 为素数}}1/p$ 发散。
\end{theorem}
\begin{proof}
若该正项级数收敛，可取 $P$ 使 $\sum_{p>P}1/p<1/2$。对任意 $X>P$，由有限乘积不等式 $\prod_i(1-u_i)\ge1-\sum_i u_i$（$0\le u_i\le1$，可归纳证明），得到
\[
 \prod_{P<p\le X}(1-1/p)>\frac12.
\]
另一方面，将有限多个收敛几何级数相乘，由唯一分解，
\[
 \sum_{n\le X}\frac1n
 \le\prod_{p\le X}(1-1/p)^{-1}
 <2\prod_{p\le P}(1-1/p)^{-1}.
\]
第一不等式成立，是因为每个 $n\le X$ 的素因子均不超过 $X$，所以其倒数必出现在乘积展开中。右端为固定常数，但调和级数发散：每个区间 $2^{j-1}<n\le2^j$ 的倒数和至少为 $1/2$。矛盾。
\end{proof}

伯特兰命题（Bertrand's postulate）进一步保证对每个整数 $n>1$，区间 $(n,2n)$ 内至少有一个素数。这里将它作为不证的选学结论，相关讨论与推广见\cite{sondow-bertrand}。它与任意长合数段并不矛盾：前者控制相对尺度，后者只要求绝对长度任意大。

\section{素数计数函数与数量级}

定义\keyterm{素数计数函数（prime-counting function）} $\pi(x)=\#\{p\le x:p\text{ 为素数}\}$（$x\ge2$）。记 $f(x)\sim g(x)$ 表示 $f(x)/g(x)\to1$。\keyterm{素数定理（prime number theorem）}给出
\begin{equation}
 \pi(x)\sim\frac{x}{\log x}\qquad(x\to\infty),
 \label{eq:prime-number-theorem}
\end{equation}
其中 $\log$ 为自然对数。该定理在此引用而不证明，进一步阅读见\cite{kedlaya-pnt}。

渐近等价描述相对误差趋于零，并不意味着两边之差趋于零，也不直接给出某个有限 $x$ 下的误差界。例如 $\pi(100)=25$，而 $100/\log100\approx21.7$。实际计数、可证明上下界与渐近式是不同层次的信息。

\begin{example}[如何正确使用渐近等价]
在允许引用素数定理的前提下，估计区间 $(x,2x]$ 中的素数个数。
\end{example}
\begin{solve}
不能把渐近符号当作精确等号直接相减，而应将每一项除以共同尺度 $x/\log x$。由素数定理，
\[
 \frac{\pi(2x)}{x/\log x}
 =\frac{\pi(2x)}{2x/\log(2x)}\frac{2\log x}{\log(2x)}\longrightarrow2,
 \qquad \frac{\pi(x)}{x/\log x}\longrightarrow1.
\]
两式相减得到 $\pi(2x)-\pi(x)\sim x/\log x$。这里两项的主导常数不同，相减后仍有非零主项；若主项完全抵消，仅凭原来的渐近等价通常不能确定差的数量级。
\end{solve}

从 $1,\ldots,N$ 等可能地选一个整数时，选到素数的概率恰为 $\pi(N)/N$，所以素数定理给出它渐近于 $1/\log N$。这里的随机性来自“等可能抽取”的实验，而不是整数本身随机地成为素数；该比例也不说明不同整数的素性事件彼此独立。

\section{习题}

本章区分无穷性、长空隙、倒数和与渐近计数四种信息。前几类结论不能自动互相推出；引用素数定理时，应明确使用的极限方向和共同尺度，而不将渐近式当作有限范围的精确公式。

\begin{enumerate}
 \item 构造连续五个合数，并验证每个数的一个真因子。
 \item 若只知道某一剩余类中有无穷多个素数，能否推出其中素数的倒数和发散？说明为何仅凭“无穷多个”不足以推断。
 \item 利用素数定理证明 $\pi(2x)-\pi(x)\sim x/\log x$；注明这里使用了一个未在正文证明的定理。
 \item 用筛法核对 $\pi(100)=25$，比较计数结果与式\eqref{eq:prime-number-theorem}的表达式。
 \item\label{ex:distribution-sparse} 给出一个无穷正整数集合，其倒数和收敛；用求和公式说明收敛，而不只写出集合名称。
 \item\label{ex:distribution-six-five} 仿照欧几里得式构造，证明模 $6$ 余 $5$ 的素数有无穷多个，不使用一般狄利克雷定理。
 \item\label{ex:distribution-density} 引用素数定理证明 $\pi(x)/x\to0$，并解释为什么这与素数倒数和发散不矛盾。
 \item\label{ex:distribution-asymptotic} 构造正函数 $f,g$ 使 $f(x)\sim g(x)$ 但 $f(x)-g(x)\to+\infty$，指出相对误差与绝对误差的区别。
\end{enumerate}

\chapter{数论算法与公钥密码入门（选学）}
\label{ch:algorithms}

前面的定理不仅回答“为什么”，也提供了计算步骤。把定理变成算法，还要规定输入范围、输出含义、何时终止，以及计算量怎样随输入增长。本章以模幂、素性检验、因数分解和公钥密码为例，展示数学正确性与计算效率之间的联系；密码部分只讨论可逐步核对的代数模型。

\section{输入规模、运算次数与算法描述}

正整数 $n$ 的二进制位数为 $L=\lfloor\log_2n\rfloor+1$，这是描述输入规模的自然参数。对最终非负的函数 $f,g$，且 $g$ 最终为正，记 $f(t)=O(g(t))$，表示存在常数 $C,t_0>0$ 使 $t\ge t_0$ 时 $f(t)\le Cg(t)$。常数不能随输入增长而改变。

试除到 $\sqrt n$ 的次数关于 $n$ 是平方根量级，却关于位数 $L$ 为指数级。欧几里得算法的非零余数每两步至少缩小一半：若本次余数超过前一除数的一半，下一次余数就小于该一半。因此对初始 $A\ge B\ge1$，除法次数为 $O(1+\log B)$。

快速模幂使用 $O(\log(e+1))$ 次模乘（$e\ge1$）。但一次大整数运算不能视为一个位操作。若模数有 $L$ 位、底数已约化，采用学校式乘法和长除法，则每次模乘可用 $O(L^2)$ 位操作实现，所以得到上界 $O(L^2\log(e+1))$。这是指定运算方法下的上界，并非最优复杂度断言。

全书基础算法可以依赖调用：扩展欧几里得算法给出最大公因数与贝祖系数；互素时系数给出模逆元；模逆元用于一次同余与中国剩余定理；快速模幂用于大指数运算。每步返回的是整数或剩余类，应明确其代表元范围。

例如，同为“循环 $n$ 次”，若 $n$ 是输入整数的数值，循环次数可能随输入位数指数增长；若 $n$ 是输入的位数，则含义完全不同。还应区分算法给出的有效输出与失败状态：模逆元算法在最大公因数大于 $1$ 时应报告不存在，不能把某个贝祖系数误当作逆元；同余合并算法应区分没有解与返回代表元 $0$，因为 $0$ 可能是一个完全合法的答案。

\section{素性检验与伪素数}

对奇数 $n>2$，取 $1<a<n$。若 $\gcd(a,n)>1$，已找到非平凡因子；若 $\gcd(a,n)=1$ 而 $a^{n-1}\not\equiv1\pmod n$，由费马小定理可判定 $n$ 合数。若同余成立，只能说通过了该底数的检验。通过检验的合数称为底数 $a$ 的\keyterm{费马伪素数（Fermat pseudoprime）}。

例如 $561=3\cdot11\cdot17$，对每个与 $561$ 互素的整数 $a$，费马小定理分别给出模 $3,11,17$ 的幂周期整除 $2,10,16$，它们均整除 $560$，故 $a^{560}\equiv1\pmod{561}$。所以仅增加若干费马检验底数不能保证正确判定素数。

\keyterm{Miller--Rabin 检验（Miller--Rabin primality test）}进一步使用平方根信息。对待测整数 $n$，先排除 $n<2$，对 $2,3$ 直接判为素数，对其他偶数判为合数。余下 $n\ge5$ 为奇数，写 $n-1=2^sd$，其中 $d$ 为奇数，取 $1<a<n-1$ 并先检验最大公因数。若互素，则计算 $b_0=a^d\pmod n$，随后令 $b_{j+1}=b_j^2\pmod n$。当 $b_0=1$，或某个 $0\le j<s$ 有 $b_j=-1\pmod n$ 时，该底数通过；其余情形判为合数。

该规则不会误判素数。若 $n$ 为素数，则 $b_s=1$；若 $b_0\ne1$，取第一次出现 $1$ 的位置，前一项是 $1$ 的平方根且不是 $1$，模素数时只能为 $-1$。合数则可能通过某些底数；重复抽样的概率分析需要额外定理，本节不以未经证明的概率界代替说明。

\begin{example}[检验失败给出了怎样的证据]
用底数 $2$ 对 $561$ 作 Miller--Rabin 检验。
\end{example}
\begin{solve}
$560=2^4\cdot35$，逐次平方得到
\[
 2^{35}\equiv263,\quad263^2\equiv166,\quad
 166^2\equiv67,\quad67^2\equiv1\pmod{561}.
\]
起点不是 $1$，而出现 $1$ 前没有出现 $560=-1$，所以判为合数。更具体地，$67$ 是模 $561$ 的一个不等于 $\pm1$ 的平方根 $1$；由 $561\mid(67-1)(67+1)$，求得 $\gcd(66,561)=33$、$\gcd(68,561)=17$，还找到了非平凡因子。本例说明额外检查平方根信息，能揭示普通费马检验遗漏的结构。
\end{solve}

\section{因数分解与算法问题的区别}

素性检验只需回答是否为素数，因数分解则要求给出非平凡因子或完整分解。两者的输出不同，不能把某个素性检验直接当作分解算法。

对奇合数 $n$，\keyterm{平方差分解（Fermat factorization）}从 $A=\lceil\sqrt n\rceil$ 开始，检查 $A^2-n$ 是否为整数平方 $B^2$；若是，则 $n=(A-B)(A+B)$，否则令 $A\leftarrow A+1$。其中 $\lceil x\rceil$ 为不小于 $x$ 的最小整数。若因子 $u,v$ 相近，则 $A=(u+v)/2$ 接近起点，因此这种方法可能较快；因子相距很远时可能很慢。例如 $6059=78^2-5^2=73\cdot83$。

进一步的 Pollard $\rho$ 方法可按如下方式试探因子：固定 $f(t)=t^2+c\pmod n$，初始化 $x=y=x_0$，反复令 $x\leftarrow f(x)$、$y\leftarrow f(f(y))$，并计算 $d=\gcd(|x-y|,n)$。若 $1<d<n$，输出因子；若 $d=n$，更换参数重试；若 $d=1$，继续迭代。实际应设置步数限额，到限额后报告此次尝试未找到因子。每个返回因子都可直接验证，但这一描述不保证任意固定参数下必定成功。

\begin{example}[观察一次因数分解迭代]
对 $n=91$ 取 $f(t)=t^2+1\pmod{91}$、$x_0=2$，执行上述迭代。
\end{example}
\begin{solve}
初始化 $x=y=2$。第一步得到 $x=5,y=26$，故 $\gcd(|5-26|,91)=\gcd(21,91)=7$，于是 $91=7\cdot13$。两个状态模 $91$ 尚不相等，但模因子 $7$ 已经相等；最大公因数把这种较小模数下的碰撞提取为因子。这个特例很快成功，并不能据此推断所有输入都只需很少几步。
\end{solve}

\section{RSA 的数论模型与正确性}

\keyterm{公钥密码（public-key cryptography）}区分公开参数与秘密参数。RSA 的基本数论模型如下：取不同奇素数 $p,q$，令 $n=pq$，选整数 $1<e<\varphi(n)$ 且 $\gcd(e,\varphi(n))=1$，再由扩展欧几里得算法求 $1\le d<\varphi(n)$ 使
\[
 ed\equiv1\pmod{\varphi(n)},\qquad \varphi(n)=(p-1)(q-1).
\]
对消息代表元 $0\le M<n$，密文 $C$ 为 $M^e$ 模 $n$ 的最小非负余数；解密输出 $C^d$ 模 $n$ 的最小非负余数。$(n,e)$ 为公开参数，$d$ 为私钥指数。

\begin{proposition}[RSA 的正确性]
\label{prop:rsa-correctness}
在上述条件下，对每个 $0\le M<n$ 均有 $M^{ed}\equiv M\pmod n$。
\end{proposition}
\begin{proof}
写 $ed=1+k(p-1)(q-1)$。模 $p$ 考察：若 $p\mid M$，两侧均为零；若 $p\nmid M$，由费马小定理有 $M^{p-1}\equiv1$，故 $M^{ed}\equiv M$。模 $q$ 同理，再由中国剩余定理合并。这一证明也覆盖与 $n$ 不互素的消息。
\end{proof}

\begin{example}
取 $p=5,q=11,e=3$，演示加密与解密。
\end{example}
\begin{solve}
$n=55,\varphi(n)=40$，由 $3\cdot27=81\equiv1\pmod{40}$，得 $d=27$。消息 $M=7$ 加密为 $C=7^3\bmod55=13$，解密 $13^{27}\bmod55=7$。若 $M=5$，虽不与 $55$ 互素，仍有 $C=15$ 和 $15^{27}\bmod55=5$，与定理一致。
\end{solve}

这里讨论的是基本代数模型，正确性不等于安全性。知道 $p,q$ 就能计算 $\varphi(n)$ 与私钥指数，这是一个明确的推导方向；不能据此未经证明地断言所有恢复消息的方法都等价于分解 $n$。实际密码方案还包含本章未讨论的消息编码与随机化机制。

\section{离散对数与综合实验}

在素数模数下固定原根 $g$，计算 $g^a\pmod p$ 可用快速模幂，而由结果求 $a\pmod{p-1}$ 是离散对数问题。Diffie--Hellman 的代数关系是：分别公开 $A=g^a\pmod p$、$B=g^b\pmod p$ 后，两方可分别计算
\[
 B^a\equiv A^b\equiv g^{ab}\pmod p.
\]
这证明两方得到相同的值，但没有证明第三方无法计算该值。讨论困难性时还需指定参数分布与计算模型，不能只根据运算方向不同作出结论。

\begin{example}[核对离散对数模型中的共同值]
取 $p=23,g=5$，两方分别选择指数 $a=6,b=15$，计算公开值与共同值。
\end{example}
\begin{solve}
$5^{11}\equiv-1$、$5^2\equiv2\pmod{23}$，所以 $5$ 是原根。公开值为 $A=5^6\bmod23=8$、$B=5^{15}\bmod23=19$。共同值为 $B^a\equiv A^b\equiv5^{90}$；由 $90\equiv2\pmod{22}$，它等于 $5^2\bmod23=2$。这些小参数便于演示运算，第三方也能直接列举指数找到答案；例题仅验证代数关系。
\end{solve}

\section{习题}

算法输出的含义要与定理的逻辑方向一致：发现费马或 Miller--Rabin 证据可以判合数，通过有限底数检验却不自动构成素性证明。RSA 与离散对数模型的计算验证说明正确性，效率和困难性仍需要相应的输入规模与计算模型。

\begin{enumerate}
 \item 写出快速模幂的逐步运算表，核查 $13^{27}\pmod{55}$。
 \item 验证 $561$ 对每个互素底数通过费马检验；找一个底数使它未通过 Miller--Rabin 检验。
 \item 用平方差法分解 $10403$，解释它为何适合此例。
 \item 实现扩展欧几里得、模逆元与中国剩余定理；测试无逆元和约束不相容的输入，并明确返回结果。
 \item\label{ex:algo-inverse-fail} 用贝祖等式解释为什么 $6$ 模 $15$ 没有逆元，并设计模逆元算法在此输入上的返回说明。
 \item\label{ex:algo-rsa-small} 对 $p=3,q=11,e=3$ 求 RSA 私钥指数 $d$，加密消息 $M=6$ 并解密，核对非互素消息也满足正确性。
 \item\label{ex:algo-sqrt-factor} 若奇合数 $n$ 满足 $r^2\equiv1\pmod n$ 且 $r\not\equiv\pm1\pmod n$，证明 $\gcd(r-1,n)$ 和 $\gcd(r+1,n)$ 都是非平凡因子。
 \item\label{ex:algo-test-cases} 为快速模幂写出覆盖负底数、零指数、非互素底数和非法模数的测试输入，说明各输入的预期结果或拒绝理由；再解释这些测试为何不能代替一般正确性证明。
\end{enumerate}

\appendix

\chapter{常用符号与结论索引}
\label{app:reference}

\section{常用符号}
\begin{description}
 \item[整数与整除] $\mathbb Z_{>0}$ 为正整数集；$a\mid b$ 表示 $b=ak$ 对某个整数 $k$ 成立；$\gcd(a,b)$ 与 $\operatorname{lcm}(a,b)$ 分别为最大公因数与最小公倍数；$v_p(n)$ 为非零整数 $n$ 的素因子 $p$ 的指数。
 \item[同余与剩余] $a\equiv b\pmod m$ 表示 $m\mid a-b$；$[a]_m$ 为剩余类；$a^{-1}\pmod m$ 仅在 $\gcd(a,m)=1$ 时存在；$\operatorname{ord}_m(a)$ 是可逆类的乘法阶。
 \item[算术函数] $\tau,\sigma,\varphi,\mu$ 分别表示约数个数、约数和、欧拉函数与莫比乌斯函数；$f*g$ 为狄利克雷卷积；$\mathbf1$ 为恒一函数，$\varepsilon$ 为卷积单位元。$\sum_{d\mid n}$ 对正约数求和。
 \item[平方剩余] $\left(\frac ap\right)$ 在奇素数分母时为勒让德符号；推广到奇正分母时为雅可比符号，符号等于 $1$ 的可解性含义随分母类型改变。
 \item[连分数] $P_k/Q_k$ 为第 $k$ 个渐近分数，下标从 $0$ 开始；$\alpha_k$ 为完整商；平方根展开中的 $B_k,C_k$ 由式\eqref{eq:surd-recurrence}定义；$\ell$ 为最小周期长度。
 \item[分析记号] $\lfloor x\rfloor$ 与 $\lceil x\rceil$ 分别为下、上取整；$\pi(x)$ 为素数计数函数；$f\sim g$ 表示比值趋于 $1$，$f=O(g)$ 表示指定极限方向下的数量级上界。
\end{description}

\section{按问题查找结论}
\begin{description}
 \item[构造整数线性组合] 用贝祖等式（定理\ref{thm:bezout}），先排除两个输入同时为零的情形。
 \item[求一次整数方程] 用定理\ref{thm:linear-diophantine}；可解条件为最大公因数整除常数项，完整解见式\eqref{eq:linear-diophantine-solutions}。
 \item[比较素因子指数] 用算术基本定理（定理\ref{thm:unique-factorization}）及式\eqref{eq:factorial-valuation}。
 \item[求一次同余与同余组] 用定理\ref{thm:linear-congruence}、\ref{thm:crt}、\ref{thm:general-crt}；记录最终模数，不能把“一个解类”理解为“一个整数”。
 \item[求约数和或反演] 用式\eqref{eq:divisor-functions}、定理\ref{thm:totient}与\ref{thm:mobius-inversion}；积性只要求互素输入相乘时成立。
 \item[化简幂或求幂同余] 用费马、欧拉定理与阶的性质；使用指标前先核查模数是否有原根，见定理\ref{thm:primitive-root-classification}。
 \item[判断平方剩余] 用定理\ref{thm:euler-criterion}和\ref{thm:quadratic-reciprocity}；勒让德符号的分母必须为奇素数。
 \item[判断两平方和] 用定理\ref{thm:two-squares}；模 $4$ 余 $3$ 的素因子指数必须为偶数。
 \item[求 Pell 方程全部解] 先按定理\ref{thm:pell-fundamental}求基本解，再按定理\ref{thm:pell-all}生成；负方程使用命题\ref{prop:negative-pell}。
\end{description}

\chapter{习题选答与综合训练}
\label{app:exercises}

\section{各章计算题选答}

\textbf{第\ref{ch:introduction}章。}相邻数问题可将 $1,\ldots,2n$ 分为 $n$ 对 $\{1,2\},\{3,4\},\ldots,\{2n-1,2n\}$ 后应用抽屉原理。证明 $\sqrt3$ 无理时，先按模 $3$ 的余数 $0,1,2$ 分类，得到平方被 $3$ 整除则原数被 $3$ 整除，再重复最小分母论证。

\textbf{第\ref{ch:divisibility}章。}$-37=8(-5)+3$；$\gcd(414,662)=2=8\cdot414-5\cdot662$。$35x+22y=1$ 的全部解为 $x=-5+22t,y=8-35t$（$t\in\mathbb Z$）。

\textbf{第\ref{ch:primes}章。}$v_2(100!)=50+25+12+6+3+1=97$；$v_5\bigl(\binom{100}{50}\bigr)=24-2\cdot12=0$。对 $1\le k<p$，$\binom pk$ 的分子含一个因子 $p$，分母 $k!(p-k)!$ 不含 $p$，故 $p\mid\binom pk$。

\textbf{第\ref{ch:congruences}章。}$17^{-1}\equiv38\pmod{43}$，因为 $17\cdot38=15\cdot43+1$。平方模 $8$ 的全部可能余数为 $0,1,4$。

\textbf{第\ref{ch:crt}章。}$14x\equiv8\pmod{30}$ 的解为 $x\equiv7,22\pmod{30}$。两条相容约束的解为 $x\equiv13\pmod{24}$，把余数 $5$ 改为 $4$ 后奇偶性冲突而无解。模 $35$ 的平方根 $1$ 有四个：$1,6,29,34$。

\textbf{第\ref{ch:arithmetic-functions}章。}$840=2^3\cdot3\cdot5\cdot7$，故 $\varphi(840)=192,\tau(840)=32,\mu(840)=0$。卷积恒等式由 $\tau=\mathbf1*\mathbf1$ 与 $\mu*\mathbf1=\varepsilon$ 得到 $\mu*\tau=\mathbf1$。

\textbf{第\ref{ch:classical-congruences}章。}$7^{222}\equiv12\pmod{13}$、$3^{100}\equiv1\pmod{20}$、$5^{117}\equiv1\pmod{19}$。若 $p>3$ 为素数，$p$ 为奇数给出 $p^2\equiv1\pmod8$，$3\nmid p$ 给出 $p^2\equiv1\pmod3$，再合并为模 $24$ 的结论。

\textbf{第\ref{ch:primitive-roots}章。}模 $7$ 的原根为 $3,5$；$10$ 模 $13$ 的阶为 $6$。$x^3\equiv1\pmod7$ 的根为 $1,2,4$；$x^2\equiv2\pmod{49}$ 的根为 $10,39$。若 $g$ 是模 $m$ 的原根，由式\eqref{eq:order-power}，$g^k$ 为原根当且仅当 $\gcd(k,\varphi(m))=1$，所以原根类有 $\varphi(\varphi(m))$ 个。

\textbf{第\ref{ch:quadratic-residues}章。}模 $13$ 的非零平方剩余为 $1,3,4,9,10,12$；$\left(\frac7{19}\right)=1$、$\left(\frac{13}{17}\right)=1$；$x^2\equiv4\pmod{45}$ 的根为 $2,7,38,43$。

\textbf{第\ref{ch:diophantine}章。}斜边为 $25$ 的三元组为 $(7,24,25)$、$(15,20,25)$。$45=3^2+6^2$、$130=3^2+11^2$，而 $75$ 不可表示，因为其素因子 $3$ 的指数为奇数。方程 $xy=2x+3y$ 化为 $(x-3)(y-2)=6$；由 $x(y-2)=3y>0$ 知 $y>2$，继而 $x>3$，所以只有正因子对，得到 $(4,8),(5,5),(6,4),(9,3)$。

\textbf{第\ref{ch:continued-fractions}章。}$415/93$ 的渐近分数为 $4,9/2,58/13,415/93$。$\sqrt2$ 的前五个渐近分数为 $1,3/2,7/5,17/12,41/29$；$\sqrt3=[1;\overline{1,2}]$、$\sqrt{13}=[3;\overline{1,1,1,1,6}]$。

\textbf{第\ref{ch:pell}章。}$D=3,5,13$ 的基本解分别为 $(2,1),(9,4),(649,180)$。负方程 $x^2-2y^2=-1$ 的前三组正解为 $(1,1),(7,5),(41,29)$。平方三角数的前两个正指标为 $n=1,8$。

\textbf{第\ref{ch:prime-distribution}章。}$6!+2,\ldots,6!+6$ 给出连续五个合数。仅由集合无穷不能推断倒数和发散，例如 $\{2^k:k\ge1\}$ 的倒数和收敛；要讨论某个素数子集，仍需关于该子集分布的额外信息。

\textbf{第\ref{ch:algorithms}章。}$13^{27}\equiv7\pmod{55}$。对 $561$ 取底数 $2$，$560=2^4\cdot35$，相继得到 $263,166,67,1$；起点非 $1$，且此前未出现 $-1\equiv560$，所以未通过 Miller--Rabin 检验。$10403=102^2-1^2=101\cdot103$。

\section{分层习题提示与方法讨论}

本节按章节给出后四题的关键步骤。计算题附结果，证明题说明主要转化；使用时应尝试补全中间推导，并回到正文检查每个所用结论的条件。习题编号由各章独立计数，以下同时注明章节。

\textbf{第\ref{ch:introduction}章。}习题\ref{ex:intro-negation}的否定为“对每个整数 $n$，存在整数 $m$ 使 $m>n$”，取 $m=n+1$ 即可；量词由“存在、任意”依次变为“任意、存在”，不等号也要同时否定。习题\ref{ex:intro-strong}先验证 $2=2\cdot1+3\cdot0$、$3=2\cdot0+3\cdot1$；对 $n\ge4$，由 $n-2\ge2$ 的表示加上一个 $2$ 得到 $n$ 的表示。两个起点分别覆盖之后的偶数与奇数。

习题\ref{ex:intro-dividing-pair}把每个所选数唯一写成 $2^ru$，其中 $u$ 为奇数。区间中可能的奇数部分只有 $1,3,\ldots,2n-1$ 这 $n$ 个；抽屉原理保证两个数具有相同的 $u$，指数较小的数就整除较大的数。习题\ref{ex:intro-descent-domain}中，正实数集合不具有正整数的良序性质；给出的值集没有最小元素。因此递降时不仅要证明参数变小，还要证明它始终属于正整数或另一个不能无限严格下降的集合。

\textbf{第\ref{ch:divisibility}章。}习题\ref{ex:div-consecutive}使用 $(n+1)-n=1$，习题\ref{ex:div-combinations}使用 $3(2n+1)-2(3n+1)=1$。任意公因数都整除右侧的 $1$，故最大公因数为 $1$；这比先分别分解两个随 $n$ 变化的整数更直接。习题\ref{ex:div-nonnegative}中，模 $4$ 得 $y$ 被 $4$ 整除，写 $y=4t,x=25-7t$，非负性给出 $t=0,1,2,3$，所以解为 $(25,0),(18,4),(11,8),(4,12)$。

习题\ref{ex:div-gcd-lcm-count}在 $d\nmid L$ 时无解。若 $d\mid L$，写 $a=du,b=dv$，则 $\gcd(u,v)=1$、$uv=L/d$。由唯一分解，$L/d$ 的每个不同素因子的整个幂次必须全部分给 $u$ 或全部分给 $v$，不能拆给两边。若不同素因子数为 $r$，有序对共有 $2^r$ 个；$L=d$ 时 $r=0$，唯一解为 $(d,d)$。

\textbf{第\ref{ch:primes}章。}习题\ref{ex:prime-cube}中 $180=2^2\cdot3^2\cdot5$，补足到 $3$ 的倍数所需的最小因子为 $2\cdot3\cdot5^2=150$，此时 $180\cdot150=30^3$。习题\ref{ex:prime-odd-divisors}把约数 $d$ 与 $n/d$ 配对；只有 $d^2=n$ 才会出现一个不与不同约数配对的对象。指数证明则使用约数个数 $\prod_i(\alpha_i+1)$，它为奇数当且仅当每个 $\alpha_i$ 都为偶数。

习题\ref{ex:prime-rational-root}设 $\sqrt n=a/b$，其中 $a,b>0$ 互素。若 $b>1$，取其素因子 $p$，由 $a^2=nb^2$ 得 $p\mid a$，与互素矛盾；所以 $b=1$。习题\ref{ex:prime-binomial}使用阶乘指数：$v_p((p^2)!)=p+1$、$v_p(p!)=1$、$v_p((p^2-p)!)=p-1$，相减得 $1$。在 $p=2$ 时计算仍成立，例如 $\binom42=6$ 恰含一个因子 $2$。

\textbf{第\ref{ch:congruences}章。}习题\ref{ex:cong-cancel}约分后应为 $x\equiv1\pmod3$，在模 $12$ 下为 $1,4,7,10$ 四类；错误保留模数只会留下 $1$。习题\ref{ex:cong-obstruction}中，平方模 $3$ 为 $0$ 或 $1$，所以平方加一不被 $3$ 整除；但 $x=2,y=1$ 已构成第二个方程的整数解，说明无解障碍依赖具体模数。

习题\ref{ex:cong-last-digit}用模 $10$ 的周期 $7,9,3,1$，因 $2026$ 除以 $4$ 余 $2$，个位为 $9$。习题\ref{ex:cong-palindrome}设数位为 $a_0,\ldots,a_{2r-1}$，回文条件为 $a_j=a_{2r-1-j}$；在交错数位和中，这两个位置的指数奇偶相反，贡献相互抵消，因此总和为零，原整数被 $11$ 整除。

\textbf{第\ref{ch:crt}章。}习题\ref{ex:crt-redundant}中 $x\equiv6\pmod8$ 已蕴含 $x\equiv2\pmod4$，所以最终解仍是 $x\equiv6\pmod8$。习题\ref{ex:crt-three}先写 $x=2+6t$，由第二式得 $6t\equiv3\pmod9$，即 $t\equiv2\pmod3$，故 $x\equiv14\pmod{18}$。再写 $x=14+18s$，模 $5$ 得 $s\equiv0\pmod5$，最终为 $x\equiv14\pmod{90}$。

习题\ref{ex:crt-zero-roots}要求 $2\mid x$ 且 $3\mid x$，即 $6\mid x$，所以模 $12$ 有 $0,6$ 两根。习题\ref{ex:crt-idempotents}在每个素数模数下有 $x(x-1)\equiv0$，因此恰有 $0,1$ 两个局部选择；中国剩余定理给出 $2^r$ 个组合。模 $30$ 的全部根为 $0,1,6,10,15,16,21,25$。

\textbf{第\ref{ch:arithmetic-functions}章。}习题\ref{ex:arith-noncoprime}的值为 $4,6,72$；因为 $12$ 与 $18$ 不互素，积性不适用，且 $72\ne4\cdot6$。习题\ref{ex:arith-invert-twelve}得到 $f(12)=F(12)-F(6)-F(4)+F(2)$，来自 $d=1,2,3,6$ 的莫比乌斯值；含平方因子的 $d=4,12$ 不贡献。习题\ref{ex:arith-squarefree}的结果为 $20-5-2=13$，对应 $1,2,3,5,6,7,10,11,13,14,15,17,19$；$d=4$ 的莫比乌斯值为零。

习题\ref{ex:arith-inverse}先在 $n=1$ 处得到 $f(1)g(1)=1$，故 $f(1)\ne0$ 为必要条件。若此条件满足，令 $g(1)=1/f(1)$，再按 $n>1$ 递推
\[
 g(n)=-\frac1{f(1)}\sum_{\substack{d\mid n\\d>1}}f(d)g(n/d).
\]
求和中 $n/d<n$，所以右侧已经确定。该定义使每个 $n>1$ 的卷积值为零，证明存在性；任何逆元都必须满足相同递推，又证明唯一性。

\textbf{第\ref{ch:classical-congruences}章。}习题\ref{ex:classical-noncoprime}中，$6^2\equiv1\pmod{35}$，故第一个余数为 $1$。模 $40$ 时拆成模 $8$ 与模 $5$：$6^{100}$ 分别余 $0$ 与 $1$，合并得到 $16$。习题\ref{ex:classical-factorial}中 $12!\equiv-1\pmod{13}$，而 $11\cdot12\equiv2$，故 $2\cdot10!\equiv-1$，乘以 $2^{-1}\equiv7$ 得 $10!\equiv6$。

习题\ref{ex:classical-fermat-all}分别证明 $2,3,5$ 整除 $a^5-a$。模 $5$ 直接用费马小定理的全整数形式；模 $3$ 若 $3\mid a$ 则两项为零，否则 $a^2\equiv1$，故 $a^5\equiv a$；模 $2$ 同理。三个模数两两互素，所以乘积 $30$ 整除差。习题\ref{ex:classical-pseudoprime}中底数 $2$ 的结论见正文；底数 $3$ 则有 $3^{340}\equiv3^{10}\equiv25\not\equiv1\pmod{31}$，因此更不可能模 $341$ 余 $1$。这里 $340\equiv10\pmod{30}$，使用的是模 $31$ 的费马周期。

\textbf{第\ref{ch:primitive-roots}章。}习题\ref{ex:order-list}对类 $1,2,4,5,7,8$ 求得的阶依次为 $1,6,3,6,3,2$，所以原根为 $2,5$。习题\ref{ex:order-composite}中，$a^k\equiv1\pmod{mn}$ 当且仅当它模 $m$、模 $n$ 都为 $1$，也就是两个局部阶都整除 $k$，最小这样的正整数即为最小公倍数。

习题\ref{ex:order-sixth}取模 $13$ 的原根 $2$，写 $x=2^t$，则 $6t\equiv0\pmod{12}$ 等价于 $t$ 为偶数，得到根 $1,3,4,9,10,12$。习题\ref{ex:order-singular}中，模 $9$ 的第一式有 $0,3,6$ 三根，第二式无根；而 $x^2\equiv0\pmod{27}$ 要求 $2v_3(x)\ge3$，等价于 $9\mid x$，所以根为 $0,9,18$。$x=0$ 单独满足方程，赋值记号本身只用于非零整数。

\textbf{第\ref{ch:quadratic-residues}章。}习题\ref{ex:quadratic-sign}有 $\left(\frac3{17}\right)=\left(\frac{17}3\right)=\left(\frac23\right)=-1$、$\left(\frac5{29}\right)=\left(\frac{29}5\right)=\left(\frac45\right)=1$；两次交换都不变号，因为至少一个素数模 $4$ 余 $1$。习题\ref{ex:quadratic-jacobi}有 $\left(\frac53\right)=\left(\frac57\right)=-1$，所以 $\left(\frac5{21}\right)=1$，但两个局部方程都无解，整体也无解。

习题\ref{ex:quadratic-two-power}可逐个检验奇数代表元，得到 $3,13,19,29$，恰有四根；也可先取 $3$，再利用取负与加 $16$ 得到另外三根。习题\ref{ex:quadratic-valuation}先由模 $3$ 条件得 $x=3y$，约去 $9$ 后要求 $y^2\equiv1\pmod3$，即 $3\nmid y$。由于 $x$ 在模 $27$ 下计数，$y$ 应在模 $9$ 下计数，共有 $1,2,4,5,7,8$ 六类；因此 $x=3,6,12,15,21,24$。

\textbf{第\ref{ch:diophantine}章。}习题\ref{ex:dio-factors}写成 $(x-y)(x+y)=45$。因 $x,y\ge0$，有 $x>y$，两个正因子应同奇且前者不大于后者；因子对为 $(1,45),(3,15),(5,9)$，得到 $(23,22),(9,6),(7,2)$。习题\ref{ex:dio-hypotenuse}中，本原斜边为奇数。若其某个素因子 $p\equiv3\pmod4$，则 $p\mid x^2+y^2$，由推论\ref{cor:prime-three-mod-four}迫使 $p\mid x,y$，与本原性矛盾。

习题\ref{ex:dio-two-square-list}中 $50=1^2+7^2$、$98=7^2+7^2$、$325=1^2+18^2$；$77=7\cdot11$ 的两个模 $4$ 余 $3$ 的素因子均为奇数指数，故不可表示。习题\ref{ex:dio-area}中，本原三角形的面积为 $uv(u^2-v^2)$。因为 $u,v$ 奇偶不同，面积被 $2$ 整除；若 $3\mid uv$ 则也被 $3$ 整除，否则 $u^2\equiv v^2\equiv1\pmod3$，差被 $3$ 整除。一般三角形的面积为本原面积乘以一个整数平方，因此仍被 $6$ 整除。

\textbf{第\ref{ch:continued-fractions}章。}习题\ref{ex:cf-negative}的答案为 $[-3;1,2]$，因为下取整是 $-3$，而不是向零截断为 $-2$。习题\ref{ex:cf-sqrt-seven}中，状态 $(B_j,C_j)$ 从 $(0,1)$ 起依次为 $(2,3),(1,2),(1,3),(2,1)$，之后回到 $(2,3)$，故 $\sqrt7=[2;\overline{1,1,1,4}]$。前四个渐近分数为 $2,3,5/2,8/3$，平方后与 $7$ 比较即可核对交错方向。

习题\ref{ex:cf-neighbors}由严格夹逼知 $cv-du$ 与 $bu-av$ 都是正整数。利用 $bc-ad=1$，有
\[
 v=b(cv-du)+d(bu-av)\ge b+d.
\]
这说明相邻行列式为 $1$ 的分数之间，要插入一个分数就必须使用较大的分母。习题\ref{ex:cf-golden}的周期尾部与原式相同，故 $\alpha=1+1/\alpha$，满足 $\alpha^2-\alpha-1=0$；由 $\alpha>1$ 得 $\alpha=(1+\sqrt5)/2$。

\textbf{第\ref{ch:pell}章。}习题\ref{ex:pell-seven}的周期长为 $4$，基本解为 $(8,3)$，第二组由 $(8+3\sqrt7)^2=127+48\sqrt7$ 得 $(127,48)$。习题\ref{ex:pell-negative-five}中最小负范数单位为 $2+\sqrt5$，正方程基本单位为 $9+4\sqrt5$，故前三组为 $(2,1),(38,17),(682,305)$；代入分别得到 $-1$。

习题\ref{ex:pell-second-order}令 $U=x_1+y_1\sqrt D$，因 $U^{-1}=x_1-y_1\sqrt D$，有 $U+U^{-1}=2x_1$，从而 $U^{n+2}-2x_1U^{n+1}+U^n=0$。比较有理部分与 $\sqrt D$ 的系数，得到两个递推；比较系数合法是因为 $\sqrt D$ 无理。习题\ref{ex:pell-obstruction}模相应素数 $p$ 得 $x^2\equiv-1\pmod p$，与 $p\equiv3\pmod4$ 矛盾。没有这样的素因子只表示这一局部障碍消失，是否有解仍须使用周期奇偶等充分判据。

\textbf{第\ref{ch:prime-distribution}章。}习题\ref{ex:distribution-sparse}可取 $\{2^j:j\ge1\}$，其倒数和为几何级数 $\sum_{j\ge1}2^{-j}=1$。习题\ref{ex:distribution-six-five}若假设素数只有有限个 $p_1,\ldots,p_r$，令 $N=6p_1\cdots p_r-1$。它与 $6$ 互素，所有素因子都模 $6$ 余 $1$ 或 $5$；若没有余 $5$ 的因子，乘积只能余 $1$。因此存在一个余 $5$ 的新素因子，与有限列表矛盾。

习题\ref{ex:distribution-density}直接得到 $\pi(x)/x\sim1/\log x\to0$。自然密度为零只控制累计个数相对于 $x$ 的比例，并不保证倒数和收敛；素数的这两个已知结论正说明两种稀疏性标准不同。习题\ref{ex:distribution-asymptotic}可取 $f(x)=x+\sqrt x$、$g(x)=x$（$x>0$），比值为 $1+1/\sqrt x\to1$，差却为 $\sqrt x\to\infty$。

\textbf{第\ref{ch:algorithms}章。}习题\ref{ex:algo-inverse-fail}若逆元存在，则 $6u-15v=1$ 有整数解，但左侧总被 $3$ 整除，因此不可能；算法应返回“逆元不存在，最大公因数为 $3$”。习题\ref{ex:algo-rsa-small}有 $n=33,\varphi(n)=20$，$3\cdot7\equiv1\pmod{20}$，故 $d=7$。消息 $6$ 加密为 $6^3\bmod33=18$，解密为 $18^7\bmod33=6$，覆盖了消息与模数不互素的情形。

习题\ref{ex:algo-sqrt-factor}若 $\gcd(r-1,n)=1$，由 $n\mid(r-1)(r+1)$ 可消去 $r-1$，推出 $r\equiv-1\pmod n$，矛盾；它又不能等于 $n$，否则 $r\equiv1$。交换两个因子得到另一个最大公因数也在 $1$ 与 $n$ 之间。习题\ref{ex:algo-test-cases}可依次取 $(a,e,m)=(-2,5,7),(5,0,13),(6,4,8),(2,3,0)$，预期输出分别为 $3,1,0$，以及因不满足 $m\ge2$ 而拒绝。有限测试能发现具体错误，而一般正确性仍来自循环不变式与终止性证明。

\section{综合例题}

\begin{example}[同余化简与解数]
求 $6x\equiv9\pmod{15}$、$x\equiv2\pmod7$ 的全部整数解，并说明模 $105$ 有几个解类。
\end{example}
\begin{solve}
第一式等价于 $x\equiv4\pmod5$。与第二式合并，得 $x\equiv9\pmod{35}$。因此模 $105$ 的解类为 $9,44,79$，共三个。模 $35$ 唯一与模 $105$ 有三类是同一整数解集的两种表达。
\end{solve}

\begin{example}[最大公因数的求和]
对 $n\ge1$，证明
\[
 \sum_{k=1}^n\gcd(k,n)=\sum_{d\mid n}d\varphi(n/d).
\]
\end{example}
\begin{proof}
按 $\gcd(k,n)=d$ 分类。写 $k=du$ 后，条件等价于 $1\le u\le n/d$、$\gcd(u,n/d)=1$，故该类含 $\varphi(n/d)$ 个数，每个贡献 $d$。对全部正约数求和即可。例如 $n=12$ 时，右侧为 $4+4+6+8+6+12=40$。
\end{proof}

\begin{example}[平方和的全部表示]
求 $x^2+y^2=2025$ 的全部非负整数解，不区分坐标次序。
\end{example}
\begin{solve}
$2025=3^4\cdot5^2$。因 $3\equiv3\pmod4$，推论\ref{cor:prime-three-mod-four}使 $x,y$ 同时被 $3$ 整除；除以 $9$ 后再次应用，得 $9\mid x,y$。写 $x=9a,y=9b$，则 $a^2+b^2=25$。规定 $a\ge b\ge0$，有 $2b^2\le25$，故只需检查 $b=0,1,2,3$，得到 $(a,b)=(5,0),(4,3)$。原方程只有 $(45,0),(36,27)$。
\end{solve}

\begin{example}[平方三角数]
求使 $n(n+1)/2$ 为平方的全部正整数 $n$ 的生成方法。
\end{example}
\begin{solve}
设 $n(n+1)/2=m^2$，配方得 $(2n+1)^2-8m^2=1$。$D=8$ 的基本解为 $(3,1)$，所以令
\[
 X_k+Y_k\sqrt8=(3+\sqrt8)^k,\qquad
 n_k=\frac{X_k-1}{2},\quad k\ge1.
\]
方程保证 $X_k$ 为奇数，故每个 $n_k$ 为正整数；反向变换和 Pell 方程的完备性保证没有遗漏。前三个指标为 $1,8,49$，对应平方三角数 $1,36,1225$。
\end{solve}

\backmatter

\begin{thebibliography}{9}
  \bibitem{mit-syllabus}
  Abhinav Kumar.
  \emph{18.781 Theory of Numbers: Syllabus}.
  MIT OpenCourseWare, Spring 2012.
  \url{https://ocw.mit.edu/courses/18-781-theory-of-numbers-spring-2012/pages/syllabus/}.

  \bibitem{mit-notes}
  Joseph Lee and Abhinav Kumar.
  \emph{18.781 Theory of Numbers: Lecture Notes}.
  MIT OpenCourseWare, Spring 2012.
  \url{https://ocw.mit.edu/courses/18-781-theory-of-numbers-spring-2012/pages/lecture-notes/}.

  \bibitem{stein-book}
  William Stein.
  \emph{Elementary Number Theory: Primes, Congruences, and Secrets}.
  Author's book page.
  \url{https://www.wstein.org/ent/}.

  \bibitem{mit-periodic}
  Joseph Lee and Abhinav Kumar.
  \emph{Periodic Continued Fractions, Quadratic Irrationalities, Lecture 20 Notes}.
  MIT OpenCourseWare, 2012.
  \url{https://ocw.mit.edu/courses/18-781-theory-of-numbers-spring-2012/resources/mit18_781s12_lec20/}.

  \bibitem{mit-four-squares}
  Joseph Lee and Abhinav Kumar.
  \emph{Four Squares Theorem, Lecture 22 Notes}.
  MIT OpenCourseWare, 2012.
  \url{https://ocw.mit.edu/courses/18-781-theory-of-numbers-spring-2012/resources/mit18_781s12_lec22/}.

  \bibitem{kedlaya-dirichlet}
  Kiran S. Kedlaya.
  \emph{Primes in Arithmetic Progressions}.
  MIT 18.785: Analytic Number Theory, 2007.
  \url{https://kskedlaya.org/18.785/dirichlet.pdf}.

  \bibitem{kedlaya-pnt}
  Kiran S. Kedlaya.
  \emph{The Prime Number Theorem}.
  MIT 18.785: Analytic Number Theory, 2007.
  \url{https://kskedlaya.org/18.785/pnt.pdf}.

  \bibitem{sondow-bertrand}
  Jonathan Sondow.
  Ramanujan Primes and Bertrand's Postulate.
  \emph{American Mathematical Monthly}, 116 (2009), 630--635.
  \url{https://arxiv.org/abs/0907.5232}.
\end{thebibliography}

\end{document}
